\RequirePackage[thinlines]{easybmat} %-- muss aufgrund eines Bugs vor etex und tikz geladen werden
\documentclass[a4paper, twoside, headsepline, index=totoc,toc=listof, fontsize=10pt, cleardoublepage=empty, headinclude, DIV=12, BCOR=5mm, titlepage]{scrartcl}
\usepackage{scrtime} % KOMA, Uhrzeit ermoeglicht
\usepackage{scrpage2} % wie fancyhdr, nur optimiert auf KOMA-Skript, leicht andere Syntax
\usepackage{etoolbox}
\usepackage{letltxmacro}
\usepackage{ifthen}

%-- Compiler-abhaengige IFs
\usepackage{ifxetex,ifluatex}
\newif\ifxetexorluatex
\ifxetex
  \xetexorluatextrue
\else
  \ifluatex
    \xetexorluatextrue
  \else
    \xetexorluatexfalse
  \fi
\fi

%--Farbdefinitionen
%-- muss vor tikz geladen werden
\usepackage[usenames, table, x11names]{xcolor}
\definecolor{dark_gray}{gray}{0.45}
\definecolor{light_gray}{gray}{0.6}

%--Zum Zeichnen
%-- muss vor polyglossia bzw. babel geladen werden
\usepackage{tikz}
\usepackage{tikz-cd}
\usetikzlibrary{external}
\tikzset{>=latex}
\usetikzlibrary{shapes,arrows,intersections}
\usetikzlibrary{calc,3d}
\usetikzlibrary{decorations.pathreplacing,decorations.markings, decorations.pathmorphing}
\usetikzlibrary{angles}

%-- Konfiguration von tikzexternalize
\tikzexternalize[prefix=tikz/, up to date check=diff]

\ifxetexorluatex
\tikzset{external/system call={xelatex \tikzexternalcheckshellescape %-- verwende LuaLaTeX, wegen dynamischer Speicherallokation
    -halt-on-error -interaction=batchmode -jobname "\image" "\texsource"}}
\else
\tikzset{external/system call={pdflatex \tikzexternalcheckshellescape 
    -halt-on-error -interaction=batchmode -jobname "\image" "\texsource"}}
\fi

%-- tikzexternalize fuer tikzcd deaktivieren, da inkompatibel
\AtBeginEnvironment{tikzcd}{\tikzexternaldisable}
\AtEndEnvironment{tikzcd}{\tikzexternalenable}

%-- um Inkompatibilitaeten von quotes und polyglossia bzw. babel zu vermeiden
\tikzset{
  every picture/.append style={
    execute at begin picture={\shorthandoff{"}},
    execute at end picture={\shorthandon{"}}
  }
}
\usetikzlibrary{quotes}
\usepackage{pgfplots}
\usepgfplotslibrary{colormaps}
\pgfplotsset{compat=1.10}



%-- Mathesymbole etc.
\usepackage{mathtools} % beinhaltet amsmath
\mathtoolsset{showonlyrefs, centercolon}
\newtagform{brackets}[\textbf]{[}{]}
\usetagform{brackets}
\usepackage{wasysym} % zusätzliche Symbole
\usepackage{amssymb} % zusätzliche Symbole
\usepackage{latexsym} % zusätzliche Symbole
\usepackage{stmaryrd} % für Blitz
\usepackage{nicefrac} % schräge Brüche
\usepackage{cancel} % Befehle zum Durchstreichen
\usepackage{mathdots} % Verbesserung von Punkten wie zB \ldots
\usepackage{mathrsfs} % das X wurde benötigt

%-- einzelne Symbole aus dem mathx-package
\DeclareFontFamily{U}{mathx}{\hyphenchar\font45}
\DeclareFontShape{U}{mathx}{m}{n}{<-> mathx10}{}
\DeclareSymbolFont{mathx}{U}{mathx}{m}{n}
\DeclareMathSymbol{\bigplus}{1}{mathx}{"90}

%-- 
\DeclareSymbolFont{bbold}{U}{bbold}{m}{n}
\DeclareSymbolFontAlphabet{\mathbbold}{bbold}
\newcommand{\mathds}[1]{\mathbb{#1}} % Um Kompatibilitaet mit frueheren Benutzung von dsfont herzustellen
\newcommand{\ind}{\mathbbold{1}} % charakteristische-Funktion-Eins

\def\mathul#1#2{\color{#1}\underline{{\color{black}#2}}\color{black}} %farbiges Untersteichen im Mathe-Modus
\renewcommand{\le}{\leqslant}
\renewcommand{\ge}{\geqslant}

%-- Underbrace als Befehl in LaTeX-Syntax (und ohne Spacing probleme mit nachfolgenden Operatoren...)
\newcommand{\Underbrace}[2]{{\underbrace{#1}_{#2}}}

%-- Alles was mit Schrift und XeTeX zu tun hat
\ifxetexorluatex
    % XeLaTeX or LuaTeX
	\ifxetex
		\usepackage{mathspec} %beinhaltet fontspec 
	\else
		\usepackage[no-math]{fontspec}
	\fi
    \usepackage{polyglossia} % moderner babel-ersatz
    \setmainlanguage[spelling=new,babelshorthands=true]{german}
	\shorthandoff{"}
	\setotherlanguage{english}
	% \newcommand\glqq{"}
	% \newcommand\grqq{"}
	\defaultfontfeatures{Mapping=tex-text, WordSpace={1.4}} %
    \setmainfont[Ligatures=Common, BoldFont={* Bold}, ItalicFont={* Light Italic},SmallCapsFont={Linux Libertine O}, SmallCapsFeatures={Letters=SmallCaps}]{Source Sans Pro}
	\setsansfont[Scale=MatchLowercase,Ligatures=Common, BoldFont={* Medium}]{Ubuntu}
	% \newfontface\fancy{texgyreadventor-italic.otf}
	% \setkomafont{title}{\fancy\Huge}
	\ifxetex
		\setallmonofonts[Scale=MatchLowercase, ItalicFont={*}]{Consolas} 
	\else
		\setmonofont[Scale=MatchLowercase, ItalicFont={*}]{Consolas}
	\fi
	\usepackage{xltxtra}
	\usepackage{fontawesome}
	\usepackage{microtype}
\else
    % default: pdfLaTeX
    \usepackage[ngerman]{babel}
    \usepackage[T1]{fontenc}
    \usepackage[utf8]{inputenc}
	\usepackage{textcomp} %verhindert ein paar Fehler bei den Fonts
    \usepackage[babel=true, tracking=true]{microtype}
	\usepackage{lmodern}
	\renewcommand{\familydefault}{\sfdefault} 
\fi


%--Mixed
\usepackage[neverdecrease]{paralist}
\usepackage[german=quotes]{csquotes}
\usepackage{makeidx}
\usepackage{booktabs}
\usepackage{wrapfig}
\usepackage{float}
\usepackage[margin=10pt, font=small, labelfont={sf, bf}, format=plain, indention=1em]{caption}
\captionsetup[wrapfigure]{name=Abb. }
\usepackage{stackrel}
\usepackage{multicol}

\flushbottom


%--Unterstreichung
\usepackage[normalem]{ulem}
\setlength{\ULdepth}{1.8pt}

%--Indexverarbeitung
\newcommand{\bet}[1]{\textbf{#1}}
\newcommand{\Index}[1]{\textbf{#1}\index{#1}}
\makeindex
\setindexpreamble{{\noindent \itshape Die \emph{Seitenzahlen} sind mit Hyperlinks zu den entsprechenden Seiten versehen, also anklickbar} \ifxetexorluatex \faHandUp\fi \par \bigskip}
\renewcommand{\indexpagestyle}{scrheadings}


%--Marginnote & todonotes
\deffootnote[1.5em]{1.5em}{1.5em}{\textsuperscript{\thefootnotemark}\ }
\usepackage[fulladjust]{marginnote}
\renewcommand*{\marginfont}{\color{dark_gray} \itshape \footnotesize}
\usepackage[textsize=small]{todonotes}
\usepackage{ragged2e}
\renewcommand*{\raggedleftmarginnote}{\RaggedLeft}
\renewcommand*{\raggedrightmarginnote}{\RaggedRight}

%--Todonotes mit tikz/externalize kompatibel machen
\LetLtxMacro{\oldtodo}{\todo}
\renewcommand{\todo}[2][]{\tikzexternaldisable\oldtodo[#1]{#2}\tikzexternalenable}
\LetLtxMacro{\oldmissingfigure}{\missingfigure}
\renewcommand{\missingfigure}[2][]{\tikzexternaldisable\oldmissingfigure[{#1}]{#2}\tikzexternalenable}

%--Konfiguration von Hyperref pdfstartview=FitH, 
\usepackage[hidelinks, pdfpagelabels,  bookmarksopen=true, bookmarksnumbered=true, linkcolor=black, urlcolor=SkyBlue2, plainpages=false,pagebackref, citecolor=black, hypertexnames=true, pdfauthor={Jannes Bantje}, pdfborderstyle={/S/U}, linkbordercolor=SkyBlue2, colorlinks=false]{hyperref}

\ifxetexorluatex
\newcommand{\hrefsym}[2]{\href{#1}{\texttt{#2\,\faExternalLink}}}
\newcommand{\hrefsymmail}[2]{\href{#1}{\texttt{\faEnvelope\,#2}}}
\renewcommand{\url}[1]{\hrefsym{#1}{\nolinkurl{#1}}}
\fi
\providecommand{\hrefsym}{\href}
\providecommand{\hrefsymmail}{\hrefsym}


% -- QR-Codes
\usepackage{qrcode} %-- hinter hyperref laden!

%--Römische Zahlen
\newcommand{\RM}[1]{\MakeUppercase{\romannumeral #1{}}}
\renewcommand{\thesection}{\Roman{section}}

%%--Abkürzungen etc.
\newcommand{\light}{\text{\Large $\lightning$}}

%-- Definitionen von weiteren Mathe-Befehlen
\DeclareMathOperator{\re}{Re} %Realteil
\DeclareMathOperator{\im}{Im} %Imaginaerteil
\DeclareMathOperator{\id}{id} %identische Abbildung
\DeclareMathOperator{\Hom}{Hom} 
\DeclareMathOperator{\Mat}{Mat}
\DeclareMathOperator{\Eig}{Eig}



\newcommand{\mathd}{\ensuremath{\mathrm{d}\mkern-1.0mu}}



%--Beweisende
\newcommand{\bewende}{\ifmmode \tag*{$\square$} \else \hfill $\square$ \fi}


%--Skalarprodukt
\DeclarePairedDelimiterX\sprod[2]{\langle}{\rangle}{#1\,\delimsize\vert\,#2}
\DeclarePairedDelimiterX\skal[2]{\langle}{\rangle}{#1\,,\,#2}

%--Betrag, Gaußklammer
\DeclarePairedDelimiter{\abs}{\lvert}{\rvert}
\DeclarePairedDelimiter{\floor}{\lfloor}{\rfloor}
\DeclarePairedDelimiter{\ceil}{\lceil}{\rceil}

%--Norm
\DeclarePairedDelimiter\doppelstrich{\Vert}{\Vert}
\newcommand{\norm}[2][\relax]{
\ifx#1\relax \ensuremath{\doppelstrich*{#2}}
\else \ensuremath{\doppelstrich*{#2}_{#1}}
\fi}


%--Umklammern
\DeclarePairedDelimiter\enbrace{(}{)}
\DeclarePairedDelimiter\benbrace{[}{]}
\DeclarePairedDelimiter\lenbrace{<}{>}
\newcommand{\ssbrace}[1]{{\scriptscriptstyle\enbrace{#1}}}

%--Mengen
\DeclarePairedDelimiterX\mengenA[1]{\lbrace}{\rbrace}{#1}
\DeclarePairedDelimiterX\mengenB[2]{\lbrace}{\rbrace}{#1\, \delimsize\vert \, #2}

\makeatletter
\newcommand{\set}[2][\relax]{
\ifx#1\relax \ensuremath{
\mengenA*{#2}}
\else \ensuremath{%
  \mengenB*{#1}{#2}}
\fi}
\makeatother

\DeclareRobustCommand{\minwidthbox}[2]{%
  \ifmmode
    \expandafter\mathmakebox
  \else
    \expandafter\makebox
  \fi
  [\ifdim#2<\width\width\else#2\fi]{#1}%
}


%--offen und abgeschlossen
\newcommand{\off}{\! \stackrel[\text{offen}]{}{\subset} \!}
\newcommand{\abg}{\! \stackrel[\text{abg.}]{}{\subset} \!}

%--Differential
\newcommand{\diff}[2]{\ensuremath{\frac{{\partial #1}}{{\partial #2}} }}
\newcommand{\diffd}[2]{\ensuremath{\frac{\mathd #1}{\mathd #2} }}

%--Diagonale Linien für Matrizen
\newcommand\tikzmark[1]{%
  \tikzexternaldisable\tikz[overlay,remember picture,baseline] \node [anchor=base] (#1) {};\tikzexternalenable}

\newcommand\MyLine[3][]{%
  \tikzexternaldisable\begin{tikzpicture}[overlay,remember picture]
    \draw[#1] (#2.south) -- (#3.north);
  \end{tikzpicture}\tikzexternalenable}
  
%--Jordankasten
\newcommand{\Jord}{ \tikzexternaldisable \begin{BMAT}{|ccc|}{|ccc|} %
				0 & & \\ %
				1\tikzmark{z} &  \diagdown & \\ %
				& \tikzmark{u}1 \MyLine[thick]{z}{u} &  0 %
			\end{BMAT} \tikzexternalenable}
			
%--Jordankasten mit Lambda
\newcommand{\JordLa}[1][\empty]{%
	\tikzexternaldisable\ifthenelse{\equal{#1}{\empty}}
		{\begin{BMAT}{|ccc|}{|ccc|} %
				\lambda & & \\ %
				1\tikzmark{z} &  \diagdown & \\ %
				& \tikzmark{u}1 \MyLine[thick]{z}{u} &  \lambda  %
			\end{BMAT}}
		{\begin{BMAT}{|ccc|}{|ccc|} %
				\lambda_{#1}  & & \\ %
				1\tikzmark{z} &  \diagdown & \\ %
				& \tikzmark{u}1 \MyLine[thick]{z}{u} &  \lambda_{#1}  %
			\end{BMAT}}\tikzexternalenable}

%--Inhaltsverzeichnis
\usepackage[tocindentauto]{tocstyle}
\usetocstyle{KOMAlike}	
\shorthandon{"}
		
%--Punkte (als letztes laden)
\usepackage{ellipsis}

\newcommand{\Semester}{SoSe 2015}
\newcommand{\fach}{Analysis 2}
\newcommand{\prof}{Prof.\ Ebert}

\providecommand{\verfasser}{Tim Keil}

%--Konfiguration von scrheadings
\setheadsepline{1pt}[\color{light_gray}]
\pagestyle{scrheadings}
\clearscrheadfoot

\providecommand{\shortFach}{\fach}
\lehead{ \includegraphics[height=0.6 cm,keepaspectratio]{../!config/Bilder/Logo_WWU_Muenster_light_gray.pdf}}
\rehead{\rule{0cm}{0.6cm}\footnotesize \sffamily \color{light_gray} \verfasser{} -- Mitschrift \shortFach}
\lohead{\rule{0cm}{0.6cm} \footnotesize \sffamily \color{light_gray} Stand: \today \; \thistime[:]}
\rohead{\includegraphics[height=0.6 cm,keepaspectratio]{../!config/Bilder/fb10logo_gray.pdf}}


\ofoot[{ \color{dark_gray} \LARGE \sffamily \thepage}]{{ \color{dark_gray} \LARGE \sffamily \thepage}} %hier wir auch der plain Stil bearbeitet!
\automark{section}
\ifoot{ \color{dark_gray} \small \leftmark}

%--Metadaten
\providecommand{\mail}{keil.menden@web.de}
\author{\verfasser}
\titlehead{\includegraphics[height=1.5cm, keepaspectratio]{../!config/Bilder/Logo_WWU_Muenster.pdf}%
\hfill \includegraphics[height=1.3cm, keepaspectratio]{../!config/Bilder/fb10logo.pdf}}
\title{Skript \fach}
\subtitle{Mitschrift der Vorlesung  \enquote{\fach} von \prof}



\begin{document}
\maketitle
\pagenumbering{Alph}
\section*{Vorwort --- Mitarbeit am Skript}
Dieses Dokument ist eine Mitschrift aus der Vorlesung \enquote{\fach, \Semester}, gelesen von \prof. Der Inhalt entspricht weitestgehend dem Tafelanschrieb. Für die
Korrektheit des Inhalts übernehme ich keinerlei Garantie! Für Bemerkungen und Korrekturen -- und seien es nur Rechtschreibfehler -- bin ich sehr dankbar. 
Korrekturen bitte durch persönliches Ansprechen oder per Mail an keil.menden@web.de.
%\begin{itemize}
%	\item Persönliches Ansprechen in der Uni oder Mails an keil.menden@web.de (gerne auch mit annotieren PDFs) \url{https://github.com/JaMeZ-B/latex-wwu}.
%	\item \emph{Direktes} Mitarbeiten am Skript: Den Quellcode poste ich auf GitHub (siehe oben), also stehen vielfältige Möglichkeiten der Zusammenarbeit zur Verfügung:
%	Zum Beispiel durch Kommentare am Code über die Website und die Kombination Fork + Pull Request. Wer sich verdient macht oder ein Skript zu einer Vorlesung, die 
%	ich nicht besuche, beisteuern will, dem gewähre ich gerne auch Schreibzugriff.
	
%	Beachten sollte man dabei, dass dazu ein Account bei \url{github.com} notwendig ist, der allerdings ohne Angabe von persönlichen Daten angelegt werden kann. 
%	Wer bei GitHub (bzw. dem zugrunde liegenden Open-Source-Programm \enquote{\texttt{git}}) -- verständlicherweise -- Hilfe beim Einstieg braucht, dem helfe ich gerne 
%	weiter. Es gibt aber auch zahlreiche empfehlenswerte Tutorials im Internet.\footnote{zB. \url{https://try.github.io/levels/1/challenges/1}, ist auf Englisch, aber dafür 
%	interaktives LearningByDoing}
%	\item \emph{Indirektes} Mitarbeiten: \TeX-Dateien per Mail verschicken. 
	
%	Dies ist nur dann sinnvoll, wenn man einen ganzen Abschnitt ändern möchte (zB. einen alternativen Beweis geben), da ich die Änderungen dann per Hand einbauen muss! Ich freue mich aber auch über solche Beiträge!
%\end{itemize}

\newpage

\tableofcontents
\cleardoubleoddemptypage
\pagenumbering{arabic}
\setcounter{page}{1}


\section{Normierte und metrische Räume} % (fold)
\label{sec:Normierte und metrische Räume}

\begin{centering}
	$\mathbb{R} \hspace{1cm} \rightarrow{} \hspace{1cm}  \mathbb{R}^n \hspace{1cm}  \rightarrow{} \hspace{1cm}  $Normierte Vektorräume$ \hspace{1cm} \rightarrow \hspace{1cm}  $Metrische Räume
\end{centering}

\subsection{Der $\mathbb{R}^n$} % (fold)
\label{sub:der_mathbb_r_n}

\[
	\mathbb{R}^n = \left\{ \left. \begin{pmatrix} x_1 \\ \vdots \\ x_n \\ \end{pmatrix}  \right| x_i \in \mathbb{R} \right\}
\]

Der $\mathbb{R}^n$ ist ein $\mathbb{R}$-Vektorraum und $\dim_{\mathbb{R}}(\mathbb{R}^n)=n$.

Sei
\[
	e_i = \begin{pmatrix} 0 \\ \vdots \\ 0 \\ 1 \\ 0 \\ \vdots \\ 0 \end{pmatrix} \leftarrow \text{$i$-ter Eintrag}
\]

\((e_1,\dots,e_n)\) ist eine Basis von $\mathbb{R}^n$ und heißt {\bfseries Standardbasis}. 

Ist \(x = \begin{pmatrix} x_1 \\ \vdots \\ x_n \end{pmatrix}\), so gilt \(x=\sum\limits^{n}_{i=1}x_ie_i\)

\subsubsection[Das Skalarprodukt]{Definition} % (fold)
\label{ssub:definition}

Das {\bfseries Skalarprodukt} auf $\mathbb{R}^n$ ist die Abbildung \[ \langle\, .\,,\,.\,\rangle : \mathbb{R}^n \times \, \mathbb{R}^n \to \mathbb{R}\] mit 

\[
 x = \left< \begin{pmatrix} x_1 \\ \vdots \\ x_n \end{pmatrix} , \begin{pmatrix} y_1 \\ \vdots \\ y_n \end{pmatrix} \right> := \sum_{i=1}^{n}x_iy_i \in \mathbb{R} 
 \]
 
es gilt:
\begin{description}
	\item[(S1)] (Bilinearität) $\forall \, x_0,x_1,y_0,y_1, \in \mathbb{R}^n \hspace{4pt} \forall \, a,b \in \mathbb{R}$
	 \begin{eqnarray}
	 \langle x_0+x_1\,,\, y_0+y_1\rangle & =& \langle x_0\,,\, y_0\rangle + \langle x_1\,,\,y_1\rangle + \langle x_1\,,\, y_0\rangle + \langle x_1\,,\, y_1\rangle \\
	 \langle ax_0\,,\, by_0\rangle &=& ab \langle x_0\,,\, y_0\rangle 
	 \end{eqnarray}
	 \item[(S2)](Symmetrie) $\forall \, x,y \in \mathbb{R}^n$ gilt $\skal{x}{y} = \skal{y}{x}$
	
	\item[(S3)](positive Definitheit) $\forall \, x \in \mathbb{R}$ gilt $\skal{x}{x} \geq 0$ und $\skal{x}{x}=0$ nur dann, wenn $x=0$
\end{description}
Desweiteren gilt $\skal{e_i}{e_j}=\delta_y := \begin{cases} 1 & i=j \\ 0 &i\neq j \end{cases}$ \hspace{1cm} (Kronecker-Symbol)

Ist \[
	x=\sum_{i=1}^{n}x_ie_i \in \mathbb{R}^n \qquad (x_i \in \mathbb{R})
	\]
so ist \begin{gather*}
	\skal{x}{e_j} = \skal{\sum_{i=1}^{n}x_ie_i}{e_j} \stackrel{(S1)}{=} \sum_{i=1}^{n}x_i \skal{e_i}{e_j} = \sum^{n}_{i=1}x_i \delta_{ij} = x_j \\
	\Rightarrow x = \sum^{n}_{j=1} \skal{x}{e_j}e_j
\end{gather*}

\subsubsection[Die $l^2$-Norm]{Definition} % (fold)
\label{ssub:definition}
Sei $x \in \mathbb{R}^n$. Die $l^2$-Norm von $x$ ist die Zahl
\[
	\norm{x}_2:= \sqrt{\skal{x}{x}} 
\]
Schreibweise manchmal auch: $\norm{x}_2 = \norm{x}$
% subsubsection definition (end)

\subsection{Euklidsche und normierte Vektorräume} % (fold)
\label{sub:euklidsche_und_normierte_vektorraume}

\subsubsection[Euklidscher Vektorraum]{Definition} % (fold)
\label{ssub:definition}
Sei $ V$  ein $\mathbb{R}$-Vektorraum. Ein {\bfseries Skalarprodukt} auf $ V$  ist eine Abbildung $\skal{.}{.}:V \times V \to \mathbb{R}$, welche die Axiome S1,S2,S3 erfüllt. \\
Ein {\bfseries euklidscher Vektorraum} ist ein Paar $(V,\skal{}{})$. $ V$  ist ein $\mathbb{R}$-Vektorraum und $\skal{}{}$ ist Skalarprodukt auf $ V$. \\
Ist $(V,\skal{}{})$ ein eukdlischer Vektorraum, so ist die durch $\skal{}{}$ induzierte Norm definiert als \[
	\norm{v}=\sqrt{\skal{v}{v}} \hspace{1cm}  v \in V 
	\]

\subsubsection[Beispiel euklidscher Vektorraum]{Beispiele} % (fold)
\label{ssub:beispiele}

\begin{itemize}
	\item $\mathbb{R}^n$ mit Standardskalarprodukt
	\item Die Abbildung \begin{align*}
		 \mathbb{C} \times \mathbb{C} &\to \mathbb{R} \\
		 (z,w) &\mapsto Re(\bar z \cdot w)
		 \end{align*}
		 ist ein Skalarprodukt. Die induzierte Norm ist der Absolutbetrag.      
\end{itemize}
% subsubsection beispiele (end)

\subsubsection[Norm]{Definition} % (fold)
\label{ssub:definition}
Sei $ V$ ein $\mathbb{R} $-Vektorraum. Eine {\bfseries Norm} auf $ V$ ist eine Abbildung $\norm{.}:V \to \mathbb{R} , v \mapsto \norm{v}$ mit
\begin{description}
	\item[(N1)]$\norm{v} \geq 0 \hspace{1cm} \forall \, v \in V$ 
	\item[(N2)] $\norm{v}=0 \Leftrightarrow v=0$
	\item[(N3)] $\norm{av}=|a|\cdot \norm{v} \hspace{1cm}  \forall\,v \in V, a \in \mathbb{R}$
	\item[(N4)] $\norm{v+w} \leq \norm{v}+\norm{w} \hspace{1cm} \forall\, v,w \in V$ \hspace{1cm} (Dreiecksungleichung) 
\end{description} 
Ein normierter Vektorraum ist ein Paar $(V,\norm{.})$ bestehend aus einem $\mathbb{R}$-VR $ V$ und einer Norm auf $ V$.

\subsubsection[Cauchy-Schwarz-Ungleichung]{Satz} % (fold)
\label{ssub:satz}
Sei $(V,\skal{}{})$ ein euklidscher Vektorraum. Für $v \in V$ setzte $\norm{v}=\sqrt{\skal{v}{v}}$ \\ Dann gilt:
\begin{enumerate}
	\item $\norm{.}$ ist eine Norm auf $ V$ 
	\item $\forall\, v,w \in V$ gilt die {\bfseries Cauchy-Schwarz-Ungleichung}:
	\[
		|\skal{v}{w}| \leq \norm{v} \cdot \norm{w}
	\]  
\end{enumerate}
\minisec{Beweis:}
\begin{enumerate}
	\item \begin{enumerate}[N1:]
		\item $\norm{v}=\sqrt{\underset{\geq 0}{\underbrace{\skal{v}{v}}}} \geq  0$
		\item $\norm{v}=\sqrt{\skal{v}{v}}=0 \Rightarrow \skal{v}{v}=0 \stackrel{\text{S3}}{\Rightarrow} v=0$ 
		\item $a \in \mathbb{R} , v \in V$ \[
			\norm{av}=\sqrt{\skal{av}{av}} \stackrel{\text{S1}}{=} \sqrt{a^2 \skal{v}{v}} = \sqrt{a^2} \sqrt{\skal{v}{v}} = \abs{a} \norm{v}
		\] 
	\end{enumerate}
	\item CS-Ungleichung ist trivial, wenn $w=0$ \\
	Ann.: $w \neq 0$: Sei $f:\mathbb{R} \to \mathbb{R} $ mit
	\begin{align*}
		f(t) &=\skal{v+tw}{v+tw} \\
		& = \skal{v}{v} +2t \skal{v}{w} + t^2 \skal{w}{w} \\
		& = \norm{v}^2+ 2t \underbrace{\skal{v}{w}}_{\in \mathbb{R}} + t^2 \norm{w}^2 
	\end{align*} 
	aus S3 folgt $f(t)\geq 0$. f ist Polynom
	\begin{align*}
		&f'(t)=0 \Leftrightarrow t_0 = - \frac{\skal{v}{w}}{\norm{w}^2} \\
		&f''(t)=2 \norm{w}^2 > 0
	\end{align*}
 $\Rightarrow$ $f$ nimmt bei $t_0$ globales Minimum an
 \begin{gather*}
 	0 \leq f(t_0)=\norm{v}^2 - \frac{2 \skal{v}{w}^2}{\norm{w}^2} + \frac{\skal{v}{w}^2}{\norm{w}^4} \norm{w}^2 = \norm{v}^2 - \frac{\skal{v}{w}^2}{\norm{w}^2} \\
	\Rightarrow \frac{\skal{v}{w}^2}{\norm{w}^2} \leq \norm{v}^2 \\ \Rightarrow \skal{v}{w}^2 \leq \norm{v}^2 \norm{w}^2 \\
	\Rightarrow \abs{\skal{v}{w}} = \sqrt{\skal{v}{w}^2} \leq \sqrt{\norm{v}^2 \norm{w}^2} = \norm{v} \cdot \norm{w}
 \end{gather*}
 \bewende
\end{enumerate}
Schließlich könnnen wir N4 zeigen:
\begin{enumerate}
	\item[N4:]\begin{align*}
		\norm{v+w}^2 &= \skal{v+w}{v+w} \\ 
		&\stackrel{\text{S1}}{=} \norm{v}^2+ 2 \skal{v}{w}+ \norm{w}^2 \\
		&\leq \norm{v} + 2 \abs{\skal{v}{w}}+ \norm{w}^2 \\
		&\stackrel{\text{CS}}{\leq} \norm{v}^2+2 \norm{v}\norm{w}+ \norm{w}^2 \\
		&=(\norm{v}+\norm{w})^2
	\end{align*} 
Wurzel ziehen zeigt N4
\bewende
\end{enumerate}
% subsubsection satz (end)

\subsubsection[Besipiele für Normen auf $\mathbb{R}$]{Weitere Beispiele für Normen auf $\mathbb{R}$} % (fold)
\label{ssub:weitere_beispiele_fur_normen_auf_mathbb_r}

\begin{gather*}
	 x = \sum^{n}_{i=1}x_ie_i \hspace{1cm}  x_i \in  \mathbb{R}
	 \end{gather*}
\begin{itemize}
	\item \text{$l^1$-Norm} $\hspace{1cm}  \norm{x}_1 = \sum^{n}_{i=1}\abs{x_i}$ 
	\item $l^{\infty}$-Norm $\hspace{1cm} \norm{x}_{\infty}=\max{\{\abs{x_i}\,|\,i=1,\dots,n}\}$
	\item Ist $p \in [1,\infty)$, so ist die $l^p$-Norm
	\[
		\norm{x}_p = \left( \sum^{n}_{i=1}\abs{x_i}^p \right)^{\frac{1}{p}} \hspace{1cm} \text{(Beweis Schwer)}
	\] 
\end{itemize}
% subsubsection weitere_beispiele_fur_normen_auf_mathbb_r (end)

\subsubsection[Äquivalenz von Normen]{Definition} % (fold)
\label{ssub:definition}
$ V$ $\mathbb{R}$-VR , $\norm{.}_0 , \norm{.}_1$ Normen auf $ V$. Die Normen $\norm{.}_0 , \norm{.}_1$ heißen {\bfseries äquivalent}, falls $C,c>0$ existieren, so dass für alle $v \in  V$ gilt:
\[
	c \norm{v}_0 \leq \norm{v}_1 \leq C \norm{v}_0 
\] 
\minisec{Bemerkung:} Das ist eine Äquivalenzrelation
\begin{itemize}
	\item {\bfseries Reflexiv}: klar
	\item {\bfseries Symmetrie}: \begin{gather*}
		c \norm{v}_0 \leq \norm{v}_1 \leq C \norm{v}_0 \\
		\Rightarrow \frac{1}{c} \norm{v}_1 \leq \norm{v}_0 \leq \frac{1}{C} \norm{v}_1
	\end{gather*}
	\item {\bfseries Transitiv}: \begin{gather*}
		c \norm{v}_0 \leq \norm{v}_1 \leq C \norm{v}_0 \hspace{0.5cm} \text{und}\hspace{0.5cm} d \norm{v}_1 \leq \norm{v}_2 \leq D \norm{v}_1  \hspace{1cm} (C,c,D,d > 0)	\end{gather*}
		\[
		\Rightarrow cd \norm{v}_0 \leq CD \norm{v}_0
		\]
\end{itemize}

\subsubsection[Beispiel Äquivalente Normen]{Beispiel} % (fold)
\label{ssub:beispiel}
Die $l^1, l^2 \, \text{und} \, l^{\infty}$-Norm auf $\mathbb{R}^n $ sind äquivalent und zwar gilt:
\[
	\norm{v}_{\infty} \leq \norm{v}_2 \leq \norm{v}_1 \leq n \norm{v}_{\infty} \hspace{1cm} \text{(Beweis Übung)}
\]
% subsubsection beispiel (end)

\subsection{Metrische Räume} % (fold)
\label{sub:metrische_raume}

\subsubsection[Abstand]{Definition} % (fold)
\label{ssub:definition}
$(V,\norm{.})$ normierter Vektorraum, $v,w \in  V$ \\
Der {\bfseries Abstand} von $v$ und $w$ ist \[
d(v,w):= \norm{v-w}	
\] 
% subsubsection definition (end)
 \subsubsection[Metrik]{Definition} % (fold)
 \label{ssub:definition}
Sei $X$ eine Menge \\
Eine {\bfseries Metrik} (oder Abstandsfunktion) ist eine Abbildung $d: X \times X \to \mathbb{R} , (x,y) \mapsto d(x,y)$, sodass gilt:
\begin{description}
	\item[(M1)] $\forall\, x,y \in X: \hspace{1cm}  d(x,y) \geq  0$
	\item[(M2)] $\forall\, x,y \in X: \hspace{1cm}  d(x,x)=0   \hspace{1cm} \text{und} \hspace{1cm} 
		 		d(x,y)=0 \Rightarrow x=y$
	\item[(M3)] $\forall\, x,y \in X: \hspace{1cm} d(x,y)=d(y,x)$
	\item[(M4)] $\forall\, x,y,z \in X: \hspace{1cm}  d(x,z) \leq d(x,y)+d(y,z) \hspace{1cm}$ ({\bfseries Dreiecksungleichung}) 
\end{description}
Ein {\bfseries Metrischer Raum} ist ein Paar $(X,d)$, $X$ Menge, $d$ Metrik auf $X$
\minisec{Beispiele:}
\begin{itemize}
	\item Ist $(V,\norm{.})$ normierter Vektorraum, so ist durch $d(v,w):= \norm{v-w}$ eine Metrik auf $ V$ gegeben.

\begin{align*}
	v,w,u \in V \hspace{1cm} d(v,u) &\stackrel{\hphantom{\text{$\triangle$-Ungl.}}}{=}  \norm{v-u} = \norm{v-w+w-u} \\
	&\stackrel{\text{$\triangle$-Ungl.}}{\leq} \norm{v-w} + \norm{w-u} \\ 
	&\stackrel{\hphantom{\text{$\triangle$-Ungl.}}}{=} d(v,w)+d(w,u)
\end{align*}
 
\minisec{Unterbeispiel:}
$V=\mathbb{R}$. Dann ist $x \mapsto \abs{x} $ eine Norm auf $\mathbb{R}$ und die induzierte Metrik auf $\mathbb{R}$ ist $d(x,y)=\norm{x-y}$
 % subsubsection definition (end)
 \item $(X,d_X)$ metrischer Raum $Y\subseteq X$ Teilmenge. Dann ist $(Y,d_y)$ ein metrischer Raum, wobei für $y_0,y_1 \in Y \text{gilt, dass} \, d_y(y_0,y_1):=d_x(y_0,y_1)$ \\
 \underline{Vereinbarung:} Ist $X \subseteq \mathbb{R}^n$ eine Teilmenge, so ist $X$ immer mit der durch das Standardskalarprodukt induzierte Metrik versehen. (es sei denn, etwas anderes ist gesagt) 
\end{itemize}

\minisec{Konvergenz}
Erinnerung (Ana I): $(x_n)_n$ Folge in $\mathbb{R}$, $x \in  \mathbb{R}$. Dann 
\[
\lim_{n \to \infty}x_n=x \, :\Leftrightarrow \, \forall\, \varepsilon > 0 \,\exists \,n_0 \, \forall\, n \geq n_0: \, \abs{x-x_1}< \varepsilon
 \] 
% subsection metrische_raume (end)
\subsubsection[Konvergenz im metrischen Raum]{Definition} % (fold)
\label{ssub:definition}
$(X,d)$ metrischer Raum, $(x_n)_n$ Folge in $X$ (d.h eine Abbildung $x: \mathbb{N} \to X$) $y \in X$. Dann Konvergiert die Folge $(x_n)_n$ gegen $y :\Leftrightarrow 
 \forall\, \varepsilon >0 \, \exists \, n_0: \forall \ n \geq n_0 \, d(x_n,y)< \varepsilon $ \\
\underline{Äquivalent}: Die Folge (in $\mathbb{R}$) $d(x_n,y)$ ist eine Nullfolge.
% subsubsection definition (end)

\subsubsection[Beispiel Metrik]{Beispiel} % (fold)
\label{ssub:3.4}
$X=\mathbb{R}^n$ mit Metrik, die durch $l^1$-$ l^2$-$ \text{ oder } l^{\infty}$- Norm induziert wird. Folge $k \mapsto x(k) \in \mathbb{R}^n$.
\[
	x(k)= \sum_{i=1}^{n}x_i(k)e_i \qquad e_i \text{ Standard-Basisvektor} \qquad x_i(k) \in  \mathbb{R} \qquad x=\sum^{n}_{i=1}x_ie_i
\]
Dann sind äquivalent
\begin{enumerate}
	\item für alle $i=1,\dots,n$ gilt $\lim\limits_{k \to \infty}x_i(k)=x_i$
	\item $\lim\limits_{k \to \infty}\norm{x(k)-x}_1=0$
	\item $\lim\limits_{k \to \infty}\norm{x(k)-x}_2=0$
	\item $\lim\limits_{k \to \infty}\norm{x(k)-x}_{\infty}=0$ 
\end{enumerate}
\minisec{Beweis:}
\underline{$1 \to 2$:} 
\[
	\norm{x(k)-x}_1:=\sum^{n}_{i=1}\underset{\stackrel{\text{(wg 1)}}{\to}0}{\underbrace{\abs{x_i(k)-x_i}}} \longrightarrow 0
\]
\\
\underline{$2 \to 3 \to 4$:} Benutze Ungleichungen $(y \in \mathbb{R}^n)$
\[
	\norm{y}_{\infty} \leq \norm{y}_2 \leq \norm{y}_1 \leq n \cdot \norm{y}_{\infty}
\]
also 
\begin{align*}
	&\lim_{k \to \infty}\norm{x(k)-x}_1=0 \\
	&\Rightarrow \lim_{k \to \infty}\norm{x(k)-x}_2=0 \\
	&\Rightarrow \lim_{k \to \infty}\norm{x(k)-x}_{\infty}=0
\end{align*}
\\
\underline{$4 \to 1$:} \begin{gather*}
	\norm{x(k)-x}_{\infty} = \max\{ \abs{x_i(k)-x} \,| i=1,\dots,n \}
	\Rightarrow \text{für alle $i$: } \abs{x_i(k)-x_i} \leq \norm{x(k)-x}_{\infty}
\end{gather*}
Also
\[
	\left[ \norm{x(k)-x}_{\infty} \stackrel{k \to \infty}{\longrightarrow}0 \right] \qquad \Rightarrow \qquad \left[ \forall\, i=1,\dots,n \quad \abs{x_i(k)-x_i} \longrightarrow 0 \right]
\]
$\Rightarrow 1)$ \bewende

\minisec{Bemerkung:} $X$ metrischer Raum $(x_n)_n , (y_n)_n$ Folgen in $X \quad \lim_{n \to \infty}x_n=x \quad \text{und} \quad \lim_{n \to \infty}y_n=y$
\[
	\Rightarrow d(x,y)=\lim_{n \to \infty}d(x_n,y_n) \qquad \text{(Blatt 2, Aufgabe 2)}
\]
Daraus folgt: \\ Ist $(x_n)_n $ Folge in $X$ und gilt $\lim_{n \to \infty}x_n=x$ sowie $\lim_{n \to \infty}x_n=x'$ dann ist $x=x'$ \\
Denn \[
	d(x,x')=\lim_{n \to \infty}\underset{=0}{\underbrace{d(x_n,x_n)}}=0
\]
\[
	\stackrel{\text{(M2)}}{\Rightarrow} x=x'
\] \bewende
% subsubsection beispiel (end)

\subsubsection[Cauchyfolge im metrischen Raum]{Definition} % (fold)
\label{ssub:definition}
Sei $(X,d)$ metrischer Raum und $(x_n)_n $ Folge in $X$. Dann heißt $(x_n)_n$ {\bfseries Cauchyfolge}, falls gilt:
\[
	\forall\, \varepsilon >0 \, \exists \, n_0 \, \forall\, n,m \geq n_0: \,\, d(x_n,x_m)<\varepsilon 
\]

\minisec{Bemerkung:} $(X,d)$ metrischer Raum $(x_n)_n $ Folgen in $X$, $\lim_{n \to \infty}x_n=x \in X$. \\
Dann ist $(x_n)_n$ Cauchyfolge, denn \\
Sei $\varepsilon >0$. Wähle $n_0$, sodass \[  \forall\, n \geq n_0 \, d(x,x_n)< \frac{\varepsilon}{2} \]
Falls $n,m \geq n_0$
\[
	d(x_n,y_n) \leq \underset{<\frac{\varepsilon}{2}}{\underbrace{d(x_n,x)}} + \underset{<\frac{\varepsilon}{2}}{\underbrace{d(x,y_n)}} < \varepsilon 
\]
% subsubsection definition (end)

\subsubsection[Vollständigkeit]{Definition} % (fold)
\label{ssub:definition}
$(X,d)$ metrischer Raum. Dann heißt $X$ {\bfseries vollständig}, falls jede Cauchy-Folge $(x_n)_n $ in $X$ gegen ein $x \in X$ konvergiert (Ana I)

\minisec{Bsp:} 

	$(0,1) \subseteq \mathbb{R}$.
	$x_n = \frac{1}{n}$ Folge in $X$ ist Cauchyfolge. \\
	$(x_n)_n$ konvergiert \underline{nicht} in $X$, denn der einzig möglicher Grenzwert (0) liegt nicht in $X$.

% subsubsection definition (end)

\subsubsection{Satz} % (fold)
\label{ssub:satz}
$\mathbb{R}^n$, mit der durch die Norm $\norm{.}_p \qquad p=1,2,\infty$ induzierten Metrik ist vollständig

\minisec{Beweis:}
\underline{$p=\infty$:} Sei $x(k)$ $l^{\infty}$-Cauchyfolge in $\mathbb{R}^n$. Dann gilt für $i=1,\dots,n$:
\[
	\abs{x_i(k)-x_i(l)} \leq \norm{x(k)-x(l)}_{\infty}
\]
Da $(x(k))_k$ $l^{\infty}$-Cauchyfolge, ist $(x_i(k))_k$ eine Cauchyfolge in $\mathbb{R}$ für alle $1\leq i \leq n$ \\
Nach dem Cauchy-Kriterium aus Analysis I sind die Folgen $(x_i(k))_k$ Konvergent. $x_i=\lim_{k \to \infty}x_i(k) \in \mathbb{R}$
\[
	x := \sum_{i=1}^{n}x_ie_i \in \mathbb{R}^n
\] 
Dann gilt 
\[
	\norm{x(k)-x}_{\infty} = \max\{ \underset{\to 0}{\underbrace{\abs{x_i(k)-x_i}}} | i=1,\dots,n \} \longrightarrow 0
\]
also
\[
	\lim_{k \to \infty}\norm{x(k)-x}_{\infty} = 0
\]
\\
\underline{$p=2$:}
$x(k)$ $l^2$-Cauchyfolge, weil $\norm{x(k)-x(l)}_{\infty} \leq \norm{x(k)-x(l)}_2$ \\
Ist $(x_k)_k$ eine $l^{\infty}$-Cauchfolge, so gibt es ein $x \in \mathbb{R}^n$ mit 
\[
	\lim_{k \to \infty}\norm{x(k)-x}_{\infty}=0
\]
\[
	\Rightarrow \norm{x(k)-x}_2 \leq n \cdot \underset{\to 0}{\underbrace{\norm{x(k)-x}_{\infty}}} \longrightarrow 0
\] 
\\
\underline{$p=1$:} Analog zum Fall $p=1$, benutze $\norm{y}_{\infty} \leq \norm{y}_1 \leq n \cdot \norm{y}_{\infty}$ 
\bewende	
% subsubsection satz (end)
\subsubsection[Stetigkeit]{Definition} % (fold)
\label{ssub:definition}
Seien $(X,d_X)$ und $(Y,d_Y)$ metrische Räume und $f:X \to Y$ eine Abbildung sowie $x \in X$.
\begin{enumerate}[(i)]
	\item $f$ heißt {\bfseries $\varepsilon$-$\delta$-stetig} (oder {\bfseries stetig}) bei $x$, wenn gilt
	\[
		\forall\, \varepsilon >0 \, \exists \, \delta>0 \, \forall\, x' \in X \text{ mit } d_X(x,x')< \delta \text{ gilt: } d_Y(f(x),f(x'))<\varepsilon 
	\]
	\item $f$ heißt {\bfseries folgenstetig} in $x$ $:\Leftrightarrow$ Ist $(x_n)_n $ eine Folge in $X$ mit $\lim_{n \to \infty}x_n=x$, so gilt
	\[
		\lim_{n \to \infty}f(x_n)=f(x)
	\]
\end{enumerate}
$f$ ist stetig (folgenstetig) auf $X$, wenn $f$ stetig (folgenstetig) bei $x$ ist für alle $x \in X$.
% subsubsection definition (end)

\subsubsection{Satz} % (fold)
\label{ssub:satz}
$f$ stetig in x $\Leftrightarrow$ $f$ folgenstetig in $x$
\minisec{Beweis} wie in Analysis I, wird bald nachgeliefert

% subsubsection satz (end)
\subsubsection{Satz} % (fold)
\label{ssub:satz}
$(X,d_X)$ , $(Y,d_Y)$, $(Z,d_Z)$ metrische Räume. 
\[ f:X \to Y \qquad g:Y \to Z \qquad x \in X \]
Falls $f$ stetig bei $x$ und stetig bei $f(x)$ so ist $g \circ f$ stetig bei $x$.
\minisec{Beweis:} Siehe Analysis I 
% subsubsection satz (end)
% subsection euklidsche_und_normierte_vektorraume (end)
\subsubsection{Beispiele} % (fold)
\label{ssub:beispiele}
\begin{align*}
&\mu: \mathbb{R}^2 \to \mathbb{R} \qquad \mu(x,y):= \, xy \qquad \text{(Multiplikation)} \\
&\alpha: \mathbb{R}^2 \to \mathbb{R} \qquad \alpha(x,y):=x+y \qquad \text{(Addition)}  \\
&q:\mathbb{R} \times \mathbb{R} \setminus {0} \to \mathbb{R} \qquad q(x,y):= \frac{x}{y} \qquad \text{(Division)}
\end{align*}
sind auf ihrem gesamten Definitionsbereich stetig.

\minisec{Beweis:}
\underline{z.\,B $q$:}
$q$ ist folgenstetig, denn: \\
Sei $(x,y) \in \mathbb{R} \times \mathbb{R} \setminus \{0\}$ (d.h $y \neq 0$) und $(x_n,y_n)$ Folge in $\mathbb{R} \times \mathbb{R} \setminus \{0\}$ mit 
\[
	\lim_{n \to \infty}(x_n,y_n)=(x,y)
\]
Dann gilt: 
\[
	\lim_{n \to \infty}x_n=x \text{ und } \lim_{n \to \infty}y_n=y \neq 0
\]
\[
	\Rightarrow q(x_n,y_n)= \frac{x_n}{y_n} \stackrel{\text{Grenzwertsätze}}{\longrightarrow } \frac{x}{y} = q(x,y)
\]
\[
	\Rightarrow q \text{ ist folgenstetig}
\]
$\alpha, \mu$ analog.
\bewende

 \minisec{\underline{Bsp:}}
 Sei $X$ metrischer Raum, $f,g:X \to \mathbb{R}$ stetig \\
 $\Rightarrow $ 
 Die Funktion 
 \[
 \begin{array}{rccl}
 	(f,g): &X &\to &\mathbb{R}^2 \\
	&x &\mapsto &(f(x),g(x))\in \mathbb{R}^2
 \end{array}
 \]
 ist stetig. \\
 Denn, $(x_n)_n $ Folge in $X$, $x_n \to x \in X$ \\ 
 $$\Rightarrow f(x_n) \to f(x) \text{ und } g(x_n) \to g(x) \qquad (f,g \text{ stetig})$$
 \[
 	\stackrel{\text{Tatsache}}{\Rightarrow } (f,g)(x_n)=(f(x_n),g(x_n)) \to (f(x),g(x)) \in \mathbb{R}^2 \\
 \]
 \[
	\Rightarrow (f,g) \text{ stetig}
 \]
 Betrachte Kompositionen 
 \[
 \begin{array}{rcccl}
 	X &\stackrel{(f,g)}{\to }  &\mathbb{R}^2 &\stackrel{\mu}{\to } &\mathbb{R} \\
	x &\stackrel{\hphantom{(f,g)}}{\mapsto } &(f(x),g(x)) &\mapsto &f(x)g(x) \\
 \end{array}
 \]
 \[
 	\Rightarrow \underset{\text{$\Rightarrow$ stetig}}{\underbrace{ \underset{\text{stetig}}{\underbrace{\mu}} \circ \underset{\text{stetig}}{\underbrace{(f,g)}} }}
	= fg
 \] 
 \[
 	\Rightarrow \text{Produkte stetiger Funktionen sind wieder stetig}
 \]
 \minisec{Analog} Summen Quotienten stetiger Funktionen sind stetig
% subsubsection beispiele (end)
\subsubsection[Lipschitz-stetig]{Definition} % (fold)
\label{ssub:definition}
$(X,d_X),(Y,d_y)$ metrische Räume $f:X \to Y$ heißt {\bfseries Lipschitz-stetig} 
$:\Leftrightarrow \exists \, L \leq 0, \text{ so dass gilt: }$
\begin{align*}
	 \forall\, x,x' \in X : \qquad 
	 d_Y(f(x),f(x')) \leq Ld_X(x,x')
\end{align*}
\minisec{Bem:} ist $f$ Lipschitz-stetig, so ist $f$ stetig: \\
sei $\varepsilon > 0$ Wähle $\delta$, so dass $\delta\,L \leq \varepsilon$ \\
Dann:
\[
	\text{Ist } d_X(x,x')<\delta \Rightarrow d_Y(f(x),f(x')) \leq L \cdot d_X(x,x') < \delta L \leq \varepsilon \qquad \Rightarrow f \text{ stetig.} 
\]
% subsubsection definition (end)

\subsubsection[Beispiel Lipschitz-stetig]{Beispiel} % (fold)
\label{ssub:beispiel}
Betrachte $\mathbb{R}^n$ mit der durch die $l^p$-Normen induzierten Metrik $p \in \{1,2,\infty\}$ \\
Dann gilt: Sind $p,q \in \{1,2,\infty\}$, so ist 
\[
	\text{id:}(\mathbb{R}^n,\norm{.}_p) \to (\mathbb{R}^n,\norm{.}_q) 
\] 
Lipschitz-stetig.
\minisec{Grund:} $\forall\, x \in \mathbb{R}^n$
\[
	\norm{x}_{\infty} \leq \norm{x}_2 \leq \norm{x}_1 \leq n \cdot \norm{x}_{\infty}
\]
% subsubsection beispiel (end)
\subsubsection{Satz} % (fold)
\label{ssub:satz}
Jede lineare Abbildung $f:\mathbb{R}^n \to \mathbb{R}^m$ ist Lipschitz-stetig.

\minisec{Beweis:} Aus LA I ist bekannt: \\
Es gibt genau ein $A \in \ \text{Mat}_{m,n}(\mathbb{R})$, so dass $f(x)=Ax \qquad (\forall\, x \in \mathbb{R}^n)$ 
\[
	\text{M}_{m,n}(\mathbb{R}) = \text{Mat}_{m,n}(\mathbb{R})
\]
% subsubsection satz (end)
\minisec{Beh:}
$\exists \, L=L_A$, so dass für alle $x \in \mathbb{R}^n$ gilt 
\[
	\norm{Ax}_2 \leq L \norm{x}_2
\]
Es gilt 
\[
	\norm{Ax}_2 \leq \norm{Ax}_1 = \sum^{m}_{i=1} \abs{(Ax)_i} = \sum^{m}_{i=1} \left|\sum^{n}_{j=1} a_{ij}x_j \right| 
\]
nach Cauchy-Schwarz gilt:
\[
	\left| \sum^{n}_{j=1}y_jz_j \right | =\abs{\skal{y}{z}} \stackrel{\text{CS}}{\leq } \norm{y}_2 \cdot \norm{z}_2 = \left( \sum^{m}_{j=1}y_j^2 \right)^{\frac{1}{2}}
	 \left( \sum^{m}_{j=1}z_j^2 \right)^{\frac{1}{2}}
\]
Also 
\[
	\Rightarrow \left| \sum^{m}_{j=1}a_{ij}x_j \right| \leq \left( \sum^{n}_{j=1}a_{ij}^2 \right)^{\frac{1}{2}} \underset{\norm{x}_2}{\underbrace{
	\left ( \sum^{n}_{j=1}x_{j}^2 \right ) }}^{\frac{1}{2}}
\]
\[
	\Rightarrow \norm{Ax}_2 \leq \sum^{n}_{i=1} \left( \left( \sum^{n}_{j=1}a_{ij}^2y \right)^{\frac{1}{2}}\norm{x}^2 \right) = \underset{=L}{\underbrace{\sum^{n}_{i=1} \left( \sum^{n}_{j=1}a_{ij}^2 \right)^{\frac{1}{2}}}} \norm{x}_2
\] 
\bewende
\subsubsection{Korollar:} % (fold)
\label{ssub:korollar}
$A \in \text{Mat}_{m,n}(\mathbb{R})$ \\
Dann ist 
\[
	\norm{A}:= \sup\limits_{\norm{x}_2 \leq 1} \norm{Ax}_2 < \infty
\]
\minisec{Beweis:}
Ist $L \geq 0$ mit $\norm{Ax}_2 \leq L \cdot \norm{x}_2 \qquad \forall\, x \in \mathbb{R}^n$ \\ 
so gilt für $\norm{x}\leq 1$: \, $\norm{Ax}_2 \leq L$
\[
	\Rightarrow \sup\limits_{\norm{x}\leq 1}\norm{Ax}_2 \leq L < \infty
\]
 \bewende
% subsubsection korollar (end)
\subsubsection[Operatornorm]{Definition} % (fold)
\label{ssub:definition}
Die Zahl $\norm{A}$ heißt {\bfseries Operatornorm} von $A$ \\
Es gilt: Die Operatornorm ist eine Norm auf dem $\mathbb{R}$-VR \, $\text{Mat}_{m,n}(\mathbb{R})$ \qquad (Übung) \\
Es gibt keine einfache Formel für $\norm{A}$
% subsubsection definition (end)

\subsection{Topologische Grundbegriffe} % (fold)
\label{sub:topologische_grundbegriffe}
\subsubsection[$\varepsilon$-Ball]{Definition} % (fold)
\label{ssub:definition}
 $(X,d)$ sei metrischer Raum, $\varepsilon >0$ und $x \in X$. Der {\bfseries $\varepsilon $-Ball} in $X$ um $x$ ist
 \[
 	B_{\varepsilon }(x) = \{ y \in X | d(x,y)< \varepsilon \}
 \]
 in $\mathbb{R}$: \[
 	B_{\varepsilon}(x)=(x-\varepsilon ,x+ \varepsilon)
 \]
% subsubsection definition (end)
\subsubsection[Innerer Punkt, Umgebung]{Definition} % (fold)
\label{ssub:definition}
$(X,D)$ metrischer Raum $Y \subseteq X, x \in X$
\begin{enumerate}
	\item $x$ heißt {\bfseries innerer Punkt} von $Y$  \[
		\exists\, \delta>0 \text{ mit } B_{\delta }(x) \subseteq Y
	\]
	In diesem Fall heißt eine {\bfseries Umgebung} von $x$ in $X$
	\item $x \in X$ heißt {\bfseries Berührpunkt} von $Y$
	\[
		\forall\, \varepsilon >0 \, \exists\, y \in Y \text{ mit } d(x,y)<\varepsilon 
	\]
	(Äquivalent: $\forall\, \varepsilon > 0 \qquad B_{\varepsilon }(x)\cap \neq \emptyset$ )
	\begin{align*}
		\mathring{Y}&= \{ x \in X \,|\, x \text{ innerer Punkt von } Y \} \\
		\overline{Y} &= \{ x \in X \,|\, x \text{ ist Berührpunkt von } Y \} 
	\end{align*}
\end{enumerate}
% subsubsection definition (end)
\subsubsection[Offen, Abgeschlossen]{Definition} % (fold)
\label{ssub:definition}
\begin{enumerate}
	\item $(X,d)$ metrischer Raum, $Y \subseteq X$. Dann heißt {\bfseries $Y$ offen} in $X$ $\Leftrightarrow $ jedes $x \in Y$ ist innerer Punkt von $Y$ \\
	( $\Leftrightarrow \forall\, x \in Y \, \exists \, \delta >0: \, B_{\delta } \subseteq Y$)
	\item $Y$ heißt {\bfseries abgeschlossen} in $X$ $:\Leftrightarrow $ Das Komplement $X-Y$ ist offen in $X$
\end{enumerate}
 \minisec{Bsp:}
 $X = \mathbb{R}$
 \begin{enumerate}
 	\item Offene Intervalle $(a,b) \subseteq \mathbb{R}$ sind offen in $\mathbb{R}$,denn
  \begin{align*}
  	&x \in (a,b) \,,\, \delta = \min \{ x-a.b-x \} >0 \\
	&\Rightarrow (x-\delta ,x+\delta ) \subseteq (a,b)
  \end{align*}  
  
 \item Abgeschlossene Intervalle $[a,b] \subseteq \mathbb{R}$ sind abgeschlossen: 
 \begin{align*}
 	&\mathbb{R} \setminus [a,b] = (-\infty,a) \cup (b, +\infty) \\
 	&\Rightarrow (x-\delta ,x+\delta ) \subseteq \mathbb{R} \setminus [a,b] \\
	&\Rightarrow x>b \text{ innerer Punkt von } \mathbb{R} \setminus[a,b] \\
	&\quad \text{ analog } x<a \text{ ist innerer Punkt von } \mathbb{R}\setminus[a,b] \\
	&\Rightarrow \mathbb{R}\setminus[a,b] \text{ offen in } \mathbb{R} \Rightarrow  [a,b] \text{ abgeschlossen in } \mathbb{R} 
\end{align*}
\end{enumerate}
% subsubsection definition (end)

\minisec{Warnung:}
Offenheit ist eine Eigenschaft von \underline{Teilmengen} metrischer Räum, nicht der Mengen selber. '$Y$ offen' sinnlos ! '$Y$ offen in $X$' sinnvoll! Analog für Abgeschlossen 
 \\
$X = \mathbb{R}^2$ \\
\begin{align*}
	&Z = \{ (x,y) \in \mathbb{R}^2\, |\, y=0 \} \\
	&Y = \{ (x,y) \in \mathbb{R}^2 \, | \, y=0,0<x<1 \} \\
	&Y \text{ offen in } Z \\
	&Y \text{ nicht offen in } X \\
	\end{align*}
% subsection topologische_grundbegriffe (end)
\subsubsection{Beispiel} % (fold)
\label{ssub:beispiel}
$X$ metrischer Raum, $ x \in X$,$\varepsilon>0$. Dann ist $B_{\varepsilon}^X(x)$ offen in $X$

\minisec{Beweis:}
zu zeigen: für jedes $y \in B_{\varepsilon}(x) \exists\, \delta >0: \, B_{\delta }(y) \subseteq B_{\varepsilon}(x)$
\[
	d(x,y)<\varepsilon \Rightarrow  \exists\, \delta>0: \, d(x,y)+\delta <\varepsilon
\] \\
Beh: $B_{\delta }(y) \subseteq B_{\varepsilon}(x)$ \\
Das ist wahr, denn
\[
	z \in B_{\delta}(y) \Rightarrow d(y,z)<\delta 
\] \\
Ferner: 
\[
	d(z,x) \leq d(z,y)+d(y,x)<\delta + d(y,x)<\varepsilon, \text{d.h.} z \in B_{\varepsilon}(x)
\]
% subsubsection beispiel (end)
\subsubsection{Satz} % (fold)
\label{ssub:satz}
Es sei $(X,d)$ metrischer Raum. \\
Dann gilt:
\begin{enumerate}
	\item $\emptyset$ und $X$ sind offen in $X$
	\item Falls $U_0,U_1$ offen in $X$ sind, so auch $U_0 \cap U_1$
	\item Falls $(U_i)_{i \in I}$ eine durch $I$ induzierte Familie von Teilmengen von $X$ ist, so ist $\bigcup\limits_{i \in I}U_i$ offen in $X$ \\
		'beliebige Vereinigungen offener Mengen sind offen'
	\item $\emptyset,X$ sind abgeschlossen in $X$
	\item $A_0,A_1$ abgeschlossen in $X$, dann ist $A_0 \cup A_1$ abgeschlossen in $X$
	\item $(A_i)_{i \in I}$, $A_i$ abgeschlossen in $X$, so ist $\bigcap\limits_{i \in I}A_i$ abgeschlossen in $X$
\end{enumerate}
\minisec{Beweis:}
\begin{enumerate}
	\item $\emptyset \subseteq X$ offen trivial \\
		$X \subseteq  X$ offen trivial
	\item Sei $x \in U_0 \cap U_1$ Dann existieren Bälle
	\begin{align*}
		B_{\varepsilon_1}^X(x) \subseteq U_1 \qquad B_{\varepsilon_0}^X(x) \subseteq U_0
	\end{align*}
Wähle nun $\varepsilon= \min \{ \varepsilon_0, \varepsilon_1 \} > 0$
\begin{align*}
	& \Rightarrow B_{\varepsilon}(x) \subseteq U_0 \cap U_1 \\
	& \Rightarrow x \text{ innerer Punkt von } U_0 \cap U_1
\end{align*}
	\item Sei $x \in \bigcup\limits_{i \in I}U_i \subseteq X$
	\[
		\Rightarrow \exists\,i \in I:\, x \in U_i, \text{ weil } U_i \subseteq X \text{ offen, gibt es } \varepsilon>0: \, B_{\varepsilon}(x) \in U_i.
	\]
Aber \[
	U_i \subseteq \bigcup\limits_{i \in I}U_i, \text{ daher: } B_{\varepsilon}(x) \subseteq \bigcup\limits_{i \in I}U_i
	\]
	\item \begin{align*}
		&\emptyset = X \setminus \underset{\text{offen nach 1)}}{\underbrace{X}} \qquad \Rightarrow \emptyset \text{ abgeschlossen in } X \\
		&X = X \setminus \underset{\text{offen}}{\underbrace{\emptyset}} \qquad \Rightarrow X \subseteq X \text{ abgeschlossen}
	\end{align*}
	\item \[
		X \setminus(A_0 \cup A_1)= \underset{\text{offen in $X$}}{\underbrace{\underset{\text{offen in $X$}}{\underbrace{(X \setminus A_0)}}\cap \underset{\text{offen in $X$}}{\underbrace{(X \setminus A_1)}}}}
	\]
	\[
		\Rightarrow A_0 \cup A_1 \subseteq X \text{ abgeschlossen}
	\]
	\item \[
		X \setminus \bigcap\limits_{i \in I}A_i = \underset{\text{offen nach 3)}}{\underbrace{\bigcup\limits_{i \in I}\underset{\text{offen in $X$}}{\underbrace{X \setminus A_i}}}}
	\]
	\[
		\Rightarrow \bigcap\limits_{i \in I}A_i \text{ abgeschlossen in } X
	\]
\end{enumerate}
% subsubsection satz (end)
\minisec{Bemerkung:}
Beliebige Durchschnitte offener Mengen sind i.A. nicht offen 
\minisec{\underline{Beispiel}:}
\[
	U_n = (-\frac{1}{n},\frac{1}{n}) \text{ ist offen in } \mathbb{R}
\]
\[
	\bigcap\limits_{n \in \mathbb{N}}U_n = \{0\} \text{ ist nicht offen in } \mathbb{R}
\]

\subsubsection{Satz} % (fold)
\label{sub:Satz}
$X$ metrischer Raum, $Y \subseteq X$. Dann sind Äquivalent:
\begin{enumerate}
	\item $Y \subseteq X$ abgeschlossen in $X$
	\item Für jede Folge $(y_n)$ mit $y_n \subseteq Y$, welche in $X$ konvergiert, gilt
	\[
		\lim_{n \to \infty}y_n \in Y
	\]
	\end{enumerate}
\minisec{Beweis:}
\minisec{Vorbemerkung:} $y \in X$ ist genau dann ein Berührpunkt von $Y \subseteq X$, wenn es eine Folge $(y_n)_n$ gibt mit $(y_n) \in Y$ und $\lim_{n \to \infty}y_n = y$
\minisec{Beweis der Vorbemerkung:}
$"\Leftarrow"$: Ist $\lim_{n \to \infty}y_n=y \qquad y_n \in Y$
\begin{align*}
	& \Rightarrow \forall\, \varepsilon>0 \,\exists\,n: \, d(y,y_n)<\varepsilon \\
	&\Rightarrow \forall\, \varepsilon>0 \, y_n \in B_{\varepsilon}(y) \cap Y \neq \emptyset \\
	&\Rightarrow y \text{ ist Berührpunkt von } Y
\end{align*}
$"\Rightarrow"$:
\begin{align*}
	&y \text{ Berührpunkt von } Y \\
	&\Rightarrow  \forall\, \varepsilon > 0 \, \exists\, y' \in Y: \, d(y,y')<\varepsilon
\end{align*}
Wähle für $n \in \mathbb{N}$ ein $y_n \in Y$ mit $d(y,y_n)<\frac{1}{n}$
\[
	\Rightarrow \lim_{n \to \infty}\underset{\in y}{\underbrace{y_n}}=y
\] \bewende
\minisec{Beweis des Satzes:}
Sei 1) nicht erfüllt. Dann:
\begin{align*}
	& \stackrel{\hphantom{\text{Vorbem.}}}{\Leftrightarrow}  Y \subseteq X \text{ nicht abgeschlossen in } X \\
	& \stackrel{\hphantom{\text{Vorbem.}}}{\Leftrightarrow} X \setminus Y \text{ nicht offen in } X \\
	& \stackrel{\hphantom{\text{Vorbem.}}}{\Leftrightarrow} \exists\,y \in X \setminus Y, \, y \text{ kein innerer Punkt von } X \setminus Y \\
	& \stackrel{\hphantom{\text{Vorbem.}}}{\Leftrightarrow} \exists\,y \in X \setminus Y, \, \forall\, \varepsilon>0 \qquad B_{\varepsilon}(y) \not\subseteq X
	 \setminus Y
	 \\
	& \stackrel{\hphantom{\text{Vorbem.}}}{\Leftrightarrow} \exists\,y \in X \setminus Y, \, \forall\, \varepsilon>0 \qquad B_{\varepsilon}(y) \cap Y \neq \emptyset
	 \\
	& \stackrel{\hphantom{\text{Vorbem.}}}{\Leftrightarrow} \exists\,y \in X \setminus Y, \, y \text{ ist Berührpunkt von } Y \\
	& \stackrel{\text{Vorbem.}}{\Leftrightarrow} \exists\, \text{ Folge } y_n \text{ in } Y, y_n \to y \notin  Y
\end{align*}
und dazu ist Äquvivalent, dass 2) nicht erfüllt ist.
\bewende
% subsection Satz (end)
\subsubsection{Satz} % (fold)
\label{ssub:satz}
$(X,d_X)$,$(Y,d_Y)$ metrischer Räume, $f: X \to Y$ Abbildung. Dann sind Äquivalent:
\begin{enumerate}
	\item $f$ ist folgenstetig.
	\item Für jede abgeschlossene Teilmenge $A \subseteq Y$ ist $f^{-1}(A) \subseteq X$ abgeschlossen.
	\item Für jede offene Teilmenge $U \in Y$ ist $f^{-1}(U)$ offen in $X$
	\item $f$ ist $\varepsilon$-$\delta$-stetig
\end{enumerate}
\minisec{Beweis:}
$\underline{1 \Rightarrow 2}$: Sei $f$ folgenstetig, $A \subseteq Y$ abgeschlossen. \\
zu zeigen: (wegen I.4.7) ist $x_n \in f^{-1}(A)$ und $x_n \to x$, so gilt $x \in f^{-1}(A)$ \\
\begin{align*}
	&x_n \to x \Rightarrow \underset{\in A}{\underbrace{f(x_n)}} \to \underset{\in y}{\underbrace{f(x)}} \qquad \text{($f$ ist folgenstetig)} \\
	&A \subseteq Y \text{ abgeschlossen } \Rightarrow  f(x) \in A \\
 	&\Rightarrow x \in f^{-1}(A) 
\end{align*}
$\underline{2 \Rightarrow 3}$:
\begin{align*}
	&U \subseteq Y \text{ offen } \Rightarrow Y \setminus U \subseteq Y \text{ abgeschlossen} \\
	&\stackrel{2)}{\Rightarrow} f^{-1}(Y \setminus U) \subseteq  X \text{ abgeschlossen.} \\
	&f^{-1}(Y \setminus U) = X \setminus f^{-1}(U) \Rightarrow f^{-1}(U) \subseteq  X \text{ offen}
\end{align*}
$\underline{3 \Rightarrow 4}$:
Sei $x \in X$, $\varepsilon>0$ 
\begin{align*}
	&\Rightarrow B_{\varepsilon}^Y(f(x)) \subseteq Y \text{ offen in } Y \\
	&\stackrel{3)}{\Rightarrow} f^{-1}(B_{\varepsilon}^Y(f(x))) \subseteq X \text{ offen} \\
	&\stackrel{\text{def. offen}}{\Rightarrow} \exists\, \delta >0 \, B_{\delta}^X(x) \subseteq f^{-1}(B_{\varepsilon}^Y(f(x)))
\end{align*}
Das heißt: ist $x' \in X$ mit $d(x,x')<\delta$
\[
	\Rightarrow d(f(x),f(x'))<\varepsilon
\]
$\underline{4 \Rightarrow 1}$:
Sei $f$ $\varepsilon$-$\delta$-stetig und $(x_n)_n$ Folge in $x$. $x_n \to x \in X$. Dann gilt $f(x_n) \to f(x)$, denn: \\
Sei $\varepsilon>0$
\begin{align*}
	&\Rightarrow \exists\,\delta >0: x' \in X \, d(x,x')< \delta \\
	&\Rightarrow d(f(x),f(x'))<\varepsilon \Rightarrow \text{weil } x_n \to x: \, \exists\,n_0: \, \forall\, n \geq n_0: \, d(x,x_n)<\delta \\
	&\Rightarrow \exists\,n_0: \, \forall\, n \geq n_0: \, d(f(x),f(x_n))<\varepsilon \text{ d.h. } f(x_n) \to f(x)
\end{align*}
% subsubsection satz (end)
\subsubsection{Beispiel} % (fold)
\label{ssub:beispiel}
$(V,\norm{.})$ normierter Vektorraum
\[
	\norm{.}: V \to \mathbb{R} \text{ stetig, denn } \norm{x}-\norm{y} \leq \norm{x-y}
\]
Also sind die Mengen (für $x \in V$, $r \geq 0$)
\[
	\{y \in V \, | \, \norm{x-y}\leq r \} \subseteq V \text{ und } \{y \in V \,|\, \norm{x-y}=r \} \subseteq V
\] abgeschlossen.
% subsubsection beispiel (end)
\subsection{Kompaktheit} % (fold)
\label{sub:kompaktheit}

\minisec{Erinnerung: Satz von Bolzano-Weierstraß}
'Jede beschränkte Folge in $\mathbb{R}$ besitzt eine konvergente Teilfolge'
\begin{itemize}
	\item $\rightarrow $ Cauchy-Konvergenzkriterium
	\item $\rightarrow $ Satz vom Minimum und Maximum (Wenn $f:A \to B$ stetig, dann gilt: $f$ nimmt Minimum und Maximum an)
	\item $\to$ Mittelwertsatz der Differentialrechnung
	\item $\to$ \underline{Satz:} $f:[a,b] \to \mathbb{R}$ stetig $\Rightarrow$ f ist gleichmäßig stetig. \\
		  $\Rightarrow$ stetige Funktionen sind Riemann-integrierbar
\end{itemize}
\subsubsection{Satz von Bolzano-Weierstraß im $\mathbb{R}^n$} % (fold)
\label{ssub:satz_von_bolzano_weierstrass_im_mathbb_r_n}
Jede beschränkte Folge in $\mathbb{R}^n$ besitzt eine konvergente Teilfolge.
 
\minisec{Beweis:}
Sei $(x(k))_k$ Folge in $\mathbb{R}^n$, $\norm{x(k)}_2 \leq C \quad  \forall\, k$. \\
$x_i(k) \quad  i=1,\dots,n$ Komponenten
\[
	\Rightarrow \underset{\text{Folge in } \mathbb{R}}{\underbrace{\abs{x_i(k)}}} \leq \norm{x(k)}^2 \leq C
\]
Nach Bolzano-Weierstraß gibt es eine Teilfolge $(x(k_l))_l$ von $x(k)$, so dass $x_1(k_l)$ (in $\mathbb{R}$) Konvergiert
Nenne \[
	x(1,l)=x(k_l)
\]
\[
\begin{matrix}
	\text{1. Komponente} & \text{2. Komponente} & \text{$n$-te Komponente} \\
	x(k_1=x(1,1) & x(2,1) & x(n,1) \\
	x(k_2)=x(1,2) & x(2,2) & x(n,2) \\
	\vdots & \vdots & \vdots 
\end{matrix}
\]
\[
	\begin{matrix}
		\text{Teilfolge} & (x(1,k))_k & (x_1(1,k)_k) \text{ Konvergiert} \\
		\text{Teilfolge} & (x(2,k))_k & (x_2(2,k)_k) \text{ Konvergiert} \\
		\vdots & \vdots & \vdots \\
		\text{Teilfolge} & (x(n,k))_k & (x_n(n,k)_k) \text{ Konvergiert}
	\end{matrix}
\]
\begin{align*}
	&\Rightarrow \text{Alle Komponentenfolgen von } (x(n,k))_k \text{ konvergieren} \\
	&\Rightarrow (x(n,k))_k \text{ ist Konvergent}
\end{align*}
\bewende
% subsubsection satz_von_bolzano_weierstrass_im_mathbb_r_n (end)
\subsubsection{Definition} % (fold)
\label{ssub:definition}
Ein metrischer Raum $X$ heißt {\bfseries folgenkompakt}, falls für jede Folge $(x_n)_n$ in $X$ eine konvergente Teilfolge hat. \\
Eine Teilmenge $Y \subseteq Y$ eines metrischen Raumes heißt folgenkompakt, falls jede Folge in $(y_n)_n$ in $Y$ eine (in $Y$!) konvergente Teilfolge besitzt.
\[
	( \Leftrightarrow Y \text{ mit der von $X$ eingeschränkten Metrik, ist folgenkompakt} )
\]
% subsubsection definition (end)
\subsubsection{Beispiel:} % (fold)
\label{ssub:beispiel}
\begin{enumerate}
	\item $[a,b]$ folgenkompakt
	\item $X= \{ x \in \mathbb{R}^n \,|\, \norm{x}_2 \leq C \}$ ist folgenkompakt: $(x_k)_k$ Folge in $X$
	\[
		\stackrel{I.5.1}{\Rightarrow } \exists\, \text{ TF } x_k \to x \in \mathbb{R}^n
	\]
	aber $X \subseteq \mathbb{R}^n$ ist abgeschlossen, also $x \in X$
	\item $(0,1)$,$(0,1]$ nicht folgenkompakt (betrachte Folge ($\frac{1}{n})_n$)
	\item $\mathbb{R}$ nicht folgenkompakt (Die Folge $(n)_n$ hat keine Konvergente TF)
	\item $\mathbb{Q} \cap [0,1] \subseteq \mathbb{R}$ nicht folgenkompakt.
	\[
		\text{Ist } x \in [0,1], \text{ so gibt es eine Folge } x_n \in [0,1] \cap \mathbb{Q} \text{ mit } x_n \to x
	\]
	Dann hat $x_n$ keine in $[0,1] \cap \mathbb{Q}$ konvergente Teilfolge (falls $x$ irrational)
\end{enumerate}
% subsubsection beispiel (end)

\subsubsection{Satz} % (fold)
\label{ssub:satz}
$X$ metrischer Raum
\begin{enumerate}
	\item Ist $X$ folgenkompakt, so ist $X$ vollständig
	\item $X$ folgenkompakt, $Y \subseteq X$ abgeschlossen $\Rightarrow$ $Y$ folgenkompakt
	\item $Y \subseteq X$, $Y$ folgenkompakt $\Rightarrow$ $Y \subseteq X$ abgeschlossen
\end{enumerate}

\minisec{Beweis:}
\begin{enumerate}
	\item Sei $(x_n)_n$ Cauchyfolge in $X$
	\[
		X \text{ folgenkompakt } \Rightarrow \text{ es gibt konvergente Teilfolge } x_{n_k} \to x \in X 
	\]
	\minisec{Beh.} $x_n \to x$ \\
	Denn \[
		d(x,x_n) \leq d(x,x_{n_k})+d(x_{n_k},x_n)
	\] Sei $\varepsilon>0$ Wähle $n_0 \quad \forall\, n,m \geq n_0: \, d(x_n,x_m)< \frac{\varepsilon}{2}$ \\
	Wähle $k: \quad n_k \geq n_0 $
	\[
		\text{für } n \geq n_k \quad d(x,x_n) \leq \underset{< \frac{\varepsilon}{2}}{\underbrace{d(x,x_{n_k})}} + 
		\underset{< \frac{\varepsilon}{2}}{\underbrace{d(x_{n_k},x)}} y \varepsilon
	\]
	\item Sei $(y_n)_n$ Folge in $Y \subseteq X$
	\[
		X \text{ folgenkompakt } \Rightarrow \exists\, \text{ TF } y_{n_k} \text{, welche gegen ein $x \in X$ konvergiert}
	\]
	Weil $Y \subseteq X$ abgeschlossen, gilt
	\[
		x = \lim_{k \to \infty}y_{n_k} \in Y \qquad \Rightarrow \qquad (y_{n_k})_k \text{ ist in $Y$ konvergent}
	\]
	$\Rightarrow Y$ folgenkompakt
	\item \underline{zu zeigen:} Ist $(y_n)_k$ Folge in $Y$. $y_n \to x \in X \Rightarrow x \in Y$ \\
	$(y_n)$ ist Cauchyfolge. \\
	 Nach 1) ist $Y$ vollständig, dh $\exists\, z \in y: \quad y_n \in z $
	\[
		\Rightarrow x=z \text{, daher } x \in Y
	\]
\end{enumerate}
% subsubsection satz (end)
\subsubsection{Satz von Heine-Borel} % (fold)
\label{ssub:satz}
Eine Teilmenge $X \subseteq \mathbb{R}^n$ ist folgenkompakt $\Leftrightarrow$ $X$ ist beschränkt und abgeschlossen in $\mathbb{R}^n$

\minisec{Beweis:}
\underline{'$\Rightarrow$'}: Sei $X$ folgenkompakt. Dann folgt aus I.5.4 3): $X \subseteq \mathbb{R}^n$ abgeschlossen \\
Wäre $X$ nicht beschränkt, dann gäbe es zu jedem $n \in \mathbb{N}$ ein $x_n \in X$ mit $\norm{x}_2 \geq n$ \\
Die Folge $(x_n)$ kann dann keine konvergente Teilfolge enthalten, also kann $X$ nicht folgenkompakt sein! \\
\\
\underline{'$\Leftarrow$'}: Sei $X \subseteq \mathbb{R}^n$ beschränkt und abgeschlossen \\
Sei $(x_n)_n$  Folge in $X$. \\
Die Folge $(x_n)_n$ ist beschränkt. 
\[
	\stackrel{I.5.1}{\Rightarrow} \text{ es gibt konvergente Teilfolge } x_{n_k} \to x \in \mathbb{R}^n
\]
Weil $X$ abgeschlossen in $\mathbb{R}^n$, ist $x \in X$
\[
	\Rightarrow X \text{ folgenkompakt}
\]
\bewende

\minisec{Beispiele:}
\begin{itemize}
	\item $D^n = \{ x \in \mathbb{R}^n \, | \, \norm{x} \leq 1 \} \subseteq \mathbb{R}^n$ \qquad \underline{abgeschlossene Einheitskugel}
	\item $S^{n-1}= \{ x \in \mathbb{R}^n \, | \, \norm{x}_2=1 \} \subseteq \mathbb{R}^n$ \qquad \underline{Einheitssphäre}
\end{itemize}
$\Rightarrow$ $D^n$,$S^{n-1}$ folgenkompakt, denn $D^n$,$S^{n-1}$ beschränkt und:
\[
	f: \mathbb{R}^n \to \mathbb{R} \qquad f(x)=\norm{x}_2 \qquad f \text{ stetig}
\]
\[
	D^n = f^{-1}(\underset{\text{abg. in $\mathbb{R}$}}{\underbrace{[0,1]}}) \qquad S^{n-1}=f^{-1}(\underset{\text{abg. in $\mathbb{R}$}}{\underbrace{\{1\}}})
\]
$\Rightarrow$ $D^n$,$S^{n-1}$ abgeschlossen
% subsubsection satz (end)

\subsubsection{Satz} % (fold)
\label{ssub:satz}
$X$,$Y$ metrische Räume, $f: X \to Y$ stetig. Sei $X$ folgenkompakt 
\[
	\Rightarrow f(X) \subseteq Y \text{ ist folgenkompakt}
\]
\minisec{Beweis:}
Sei $y_n$ Folge in $f(X)$ \\
\[ \forall\, n \, \exists\, x_n \in X \text{ mit } f(x_n)=y_n \] 
$(x_n)_n$ Folge in $X$
\[
	X \text{ folgenkompakt } \Rightarrow  \exists\, \text{ konvergente Teilfolge } x_{n_k} \to x \in X
\]
\[
	f \text{ stetig } \Rightarrow y_{n_k} = f(x_{n_k}) \to f(x) \in f(X)
\]
Damit folgt die Behauptung
\bewende

% subsubsection satz (end)
\subsubsection{Satz vom Minimum und Maximum} % (fold)
\label{ssub:satz_vom_minimum_und_maximum}
$X$ folgenkompakt, $f: X \to \mathbb{R}$ stetig. \\ Dann existieren $x_{\min},x_{\max} \in X$, so dass für alle $x \in X$ gilt:
\[
	f(x_{\min}) \leq f(x) \leq x_{\max}
\]
\minisec{Beweis:}
$f(X) \subseteq \mathbb{R}$ ist folgenkompakt (Satz I.5.6) \\
Nach I.5.5 ist $f(X) \subseteq \mathbb{R}$ abgeschlossen und beschränkt. \\
Seien $c,C \in \mathbb{R}$ definiert durch
\begin{align*}
	c &:= \inf f(x) \\
	C &:= \sup f(x)
\end{align*}
Dann sind nach Ana I $c,C$ Berührpunkte um $f(x)$
\[
	f(x) \text{ abgeschlossen } \Rightarrow c,C \in f(x)
\]
Wähle nun $x_{\min},x_{\max} \in X$ mit
\[
	f(x_{\min})=c \,,\,f(x_{\max})=C
\] \bewende
% subsubsection satz_vom_minimum_und_maximum (end)

\subsubsection{Satz} % (fold)
\label{ssub:satz}
Je zwei Normen auf $\mathbb{R}^n$ sind äquivalent

\minisec{Beweis:}
Da die Äquivalenz von Normen die Äquivalenzrelation ist, genügt es zu zeigen, dass jede Norm auf $\mathbb{R}^n$ zur $l^2$-Norm äquivalent ist. \\
Sei $\norm{.}:\mathbb{R}^n \to  \mathbb{R}$ diese Norm. \\
zu zeigen: $\exists\, c,C>0 \forall\, x \in \mathbb{R}^n: \qquad c \norm{x}_2 \leq \norm{x} \leq C \norm{x}_2$ \\
Sei \[
	x = \sum^{n}_{i=1}x_ie_i \qquad , x_i \in \mathbb{R}
\]
Dann gilt:
\[
	\norm{x}= \norm{\sum^{n}_{i=1}x_ie_i} \leq \sum^{n}_{i=1}\norm{x_ie_i} = \sum^{n}_{i=1}\abs{x_i} \, \norm{e_i} 
	\leq \underset{\norm{x}_2}{\underbrace{\left(\sum^{n}_{i=1}\abs{x_i}^2  \right)^{\frac{1}{2}}}} \cdot
	 \underset{=:C >0}{\underbrace{\left(\sum^{n}_{i=1}\norm{e_i}^2  \right)^{\frac{1}{2}}}} = C \norm{x}_2
\]
Weiter ist 
\begin{align*}
	\norm{.}:\mathbb{R}^n &\to  \mathbb{R} \\
	x &\mapsto \norm{x}
\end{align*}
ist stetig, denn $(\norm{x}-\norm{y}) \leq \norm{x-y} \leq C \norm{x-y}_2$ \\
Nach dem Satz vom Minimum und Maximum gibt es $x_{\min} \in S^{n-1}$ mit: 
\[
	\forall\, x \in S^{n-1}: \norm{x} \geq \norm{x_{\min}} =: c>0 \text{ (da $x_{\min}\neq 0$)}
\]
Sei $y \in \mathbb{R}^n$ beliebig aber $y \neq 0$. Dann folgt:
\[
	\norm{y}= \bigg\vert\bigg\vert \underset{\in \mathbb{R}}{\underbrace{\norm{y}_2}} \underset{\in S^{n-1}}{\underbrace{\frac{y}{\norm{y}_2}}} \bigg\vert\bigg\vert  = \norm{y}_2 \cdot
	 \norm{\frac{y}{\norm{y}_2}} \geq \norm{y}_2 \cdot \norm{x_{\min}} = c \cdot \norm{y}_2
	\]
	\bewende
% subsubsection satz (end)


\subsubsection{Erinnerung/Definition} % (fold)
\label{ssub:erinnerung_definition}

% subsubsection erinnerung_definition (end)
$X$ metrischer Raum, $f:X \to \mathbb{R}$ heißt {\bfseries gleichmäßig stetig}, wenn gilt:
\[
	\forall\, \varepsilon>0 \exists\, \delta >0 \,: \forall\, x,x' \in X: \, \norm{x-x'}<\delta \Rightarrow \norm{f(x)-f(x')}<\varepsilon
\]
\minisec{Behauptung:}
$X$ folgenkompakt, $f:X \to \mathbb{R}$ stetig $\Rightarrow $ $f$ ist gleichmäßig stetig
\begin{gather*}
	X \stackrel{f}{\to} \mathbb{R} \text{ stetig, } \varepsilon>0 \\
	x \in X: \qquad U_x = \{x' \in X \,|\,\abs{f(x')-f(x)}< \frac{\varepsilon}{2} \} = f^{-1}((f(x)-\frac{\varepsilon}{2},f(x)+\frac{\varepsilon}{2})) 
	\stackrel{\text{offen, $f$ stetig}}{\subseteq } X
\end{gather*}
\underline{Beispiel:} $f(x)=\frac{1}{x}$ \\
Um zu beweisen, dass $f$ glm stetig ist, brauchen wir:
\[
	\exists\, \delta >0 \,\forall\, y \in X \text{ gibt es } x \in X: \qquad B_{\delta}(y)\subseteq U_x 
\]
Denn 
\begin{align*}
	d(y,y')<\delta &\Rightarrow y' \in B_{\delta}(y) \subseteq U_x \\
					& \Rightarrow  \abs{f(y')-f(y)} \leq \underset{<\frac{\varepsilon}{2}}{\underbrace{\abs{f(y')-f(x)}}} +
					 \underset{<\frac{\varepsilon}{2}}{\underbrace{\abs{f(x)-f(y)}}} < \varepsilon
\end{align*}
 \subsubsection{Definition} % (fold)
 \label{ssub:definition}
 $X$ metrischer Raum, Eine {\bfseries offene Überdeckung} von $X$ ist eine Familie $U=(U_i)_{i \in I}$, $U_i \subseteq X$ offen für alle $i \in I$ und es gilt:
 \[
 	\bigcup\limits_{i \in I}U_i = X
 \]
 \minisec{Beispiel:}
 \begin{align*}
 	&f:X \to \mathbb{R} \text{ stetig, } x \in X, \,\varepsilon>0 \\
	&U_x = \{ x' \in X \,|\,\abs{f(x)-f(x')}< \frac{\varepsilon}{2} \} \subseteq X \\
	&(U_x)_{x \in X} = U \text{ offene Überdeckung von }X
 \end{align*}
 % subsubsection definition (end)
 \subsubsection{Satz Lebesgue-Lemma} % (fold)
 \label{ssub:satz_lebesgue_lemma}
 Sei $X$ ein folgenkompakter, metrischer Raum und $U=(U_i)_{i \in I}$ offene Überdeckung von $X$. Dann existiert $\delta >0$, so dass gilt:
 \[
 	\forall\, x \in X \, \exists\,i \in I: \qquad B_{\delta}(x)\subseteq U_i
 \]
 Ein solches $\delta$ heißt {\bfseries Lebesgue Zahl} der Überdeckung von $U$
 % subsubsection satz_lebesgue_lemma (end)
 
 
 \minisec{Beweis:}
 Durch Wiederspruch: \\
 \underline{Annahme:} $X$ folgenkompakt, $U=(U_i)_{i \in I}$ offene Überdeckung, welche keine Lebesgue-Zahl hat \\
 $\Rightarrow \forall\, \delta >0 \exists\, x \in X:$ $B_{\delta}(x)$ ist in keiner der Mengen $U_i$ enthalten \\
 $\forall\, n \in \mathbb{N} \, \exists\, x_n \in X:$ $B_{\frac{1}{n}}(x_n)$ ist in keinem $U_i, i \in I$ enthalten. \\
 Da $X$ folgenkompakt ist, gibt es eine konvergente Teilfolge $(x_{n_k})_k$ so dass $B_{\frac{1}{n_k}}(x_{n_k})$ in keinem $U_i$ enthalten ist.
 \begin{align*}
 	x = \lim_{k \to \infty}x_{n_k} \in X & \Rightarrow \exists\,i_0 \in I \qquad ,x \in U_{i_0} \subseteq X \text{ offen}
	& \exists\,\varepsilon>0 \qquad B_{\varepsilon}(x) \subseteq U_{i_0}
 \end{align*}  
 Das ist ein Widerspruch, denn \\
 Wähle h so, dass $\frac{1}{n_k}< \varepsilon$ und $d(x,x_{n_k}) < \frac{\varepsilon}{2}$ \\
 $\Rightarrow B_{\frac{1}{n_k}}(x_{n_k}) \subseteq B_{\varepsilon}(x)=U_{i_0}$
 \[
 	d(y,x_{n_k})<\frac{1}{n_k} \Rightarrow d(y,x) \leq d(y,x_{n_k})+d(x_{n_k},x) < \frac{1}{n_k}+\frac{\varepsilon}{2} < \frac{\varepsilon}{2}+ \frac{\varepsilon}{2} = \varepsilon \,  \lightning
 \] 
 \bewende
 
 \subsubsection{Korollar} % (fold)
 \label{ssub:korollar}
 $X$ folgenkompakt
 \[
 	f:X \to \mathbb{R} \text{ stetig } \Rightarrow f \text{ gleichmäßig stetig} 
 \]
 \minisec{Beweis:} s.o
 % subsubsection korollar (end)
 
\subsubsection{Definition} % (fold)
\label{ssub:definition}
$X$ metrischer Raum, $U=(U_i)_{i \in I}$ eine offene Überdeckung. Eine {\bfseries Teilüberdeckung von U} ist eine offene Überdeckung $V=(V_j)_{j \in J}$ von $X$, so dass gilt: 
\[
	\forall\, j \in J \,\exists\,i \in I: V_j = U_i
\]
% subsubsection definition (end)
\subsubsection{Definition} % (fold)
\label{ssub:definition}
Ein metrischer Raum heißt {\bfseries überdeckungskompakt} (oder {\bfseries Kompakt}) falls {\bfseries jede} offene Überdeckung $U=(U_i)_{i \in I}$ eine endliche Teilüberdeckung $V=(V_j)_{j \in J}$ besitzt. (d.h die Indexmenge J ist endlich)
% subsubsection definition (end)
\subsubsection{Satz von Heine-Borel} % (fold)
\label{ssub:satz_von_heine_borel}
$X$ metrischer Raum. Dann sind äquivalent:
\begin{enumerate}
	\item $X$ folgenkompakt
	\item $X$ ist überdeckungskompakt
\end{enumerate}
\minisec{Beweis:} 
\underline{$2 \Rightarrow 1$}:
Wir zeigen nun:
\[
	\text{überdeckungskompakt } \Rightarrow \text{ folgenkompakt}
\]
Durch Widerspruch: Sei $X$ überdeckungskompakt, und $(x_n)_n$ Folge in $X$, welche keine konvergente Teilfolge hat. \\
\underline{Beobachte:}
$(x_n)_n$ hat keine Konvergente Teilfolge
\[
	\Rightarrow \forall\, x \in X \,\exists\, \varepsilon>0 \text{, so dass }x_n \in B_{\varepsilon}(x) \text{ nur für endlich viele $n$ gilt}
\]
\[
	\Rightarrow \forall\, x \in X \,\exists\, \varepsilon_x>0 \text{ ,so dass }x_n \in B_{\varepsilon}(x) \text{ nur für unendliche viele $n \in \mathbb{N}$ golt.}
\]
Dann ist:
\[
	U = (B_{\varepsilon_x}(x))_{x \in X} 
\]
eine offene Überdeckung von $X$ \\
Weil $X$ Überdeckungskompakt ist, gibt es endliche Teilüberdeckung
\[
	\Rightarrow  \exists\,y_1,\dots,y_r \in X \text{ mit:}
\]
\[
	X = B_{\varepsilon_{y_1}}(y_1) \cup \dots \cup B_{\varepsilon_{y_r}}(y_r)
\]
\bewende
\minisec{Bemerkung:}
$X$ metrischer Raum. Man kann zeigen, dass folgendes äquvivalent ist:
\begin{enumerate}
	\item $X$ kompakt
	\item $X$ ist vollständig und total beschränkt
	\item $X$ hat folgende Eigenschaft:
	\[
		\text{Sind }   A_1  \supseteq A_2 \supseteq A_3 \supseteq \dots \text{ abgeschlossene Teilmengen von $X$ und }A_n \neq 0 \, \forall\, n
		\]
		\[
			\bigcap\limits_{n=1}^{\infty}A_n \neq \emptyset
		\]
\end{enumerate}
ohne Beweis \\
\\
\underline{$1 \Rightarrow 2$}: Wir setzen hierfür Lemma I.5.16 voraus.
$X$ folgenkompakt, $U=(U_i)_{i \in I}$ offene Überdeckung von $X$ \\
\underline{zu zeigen:} Es gibt endliche Teilüberdeckung $V=(V_j)_{j \in J}$ \\
Sei $\delta>0$ eine Lebesgue-Zahl von $U$ (existiert nach I.5.11)
\[
	5.16 \Rightarrow \exists\, x_1, \dots, x_r \in X \qquad \bigcup\limits_{j=1}^r B_{\delta}(x_j)=X
\]
\[
	\Rightarrow \forall\, j=1,\dots,n \, \exists\, i,j \in I: \qquad B_{\delta}(x_j) \subseteq U_{ij}
\]
\[
	\Rightarrow X = U_{i_1} \cup \cdots \cup U_{i_1}
\]
\bewende

\subsubsection{Lemma} % (fold)
\label{ssub:lemma}
$X$ folgenkompakter metrischer Raum. \\
Dann ist $X$ {\bfseries total beschränkt}, d.h. $\forall\, \varepsilon>0$ gibt es endlich viele $x_1,\dots,x_r \in X$, so dass:
\[
	\bigcup\limits_{i=1}^r B_{\varepsilon}(x_i)=X
\]
\minisec{Beweis:} Wir zeigen: 
\[
	X \text{ nicht total beschränkt } \Rightarrow X \text{ nicht folgenkompakt}
\]
Sei $X$ also nicht total beschränkt. $\Rightarrow \exists\,>0 \text{, so dass gilt:}$
\begin{align*}
		&\text{Sind } x_1,\dots,x_r \in X \text{ beliebig} \\
		&\Rightarrow \bigcup\limits_{i=1}^rB_{\varepsilon}(x_i) \neq X \\
		&\Rightarrow X \setminus \bigcup\limits_{i=1}^rB_{\varepsilon}(x_i) \neq \emptyset
\end{align*}
Wähle $x_1 \in X \Rightarrow X \setminus B_{\varepsilon}(x_1) \neq \emptyset$ \\
Wähle $x_2 \in X \setminus B_{\varepsilon(x_1)} \Rightarrow X \setminus (B_{\varepsilon}(x_1)\cup(B_{\varepsilon}(x_2))) \neq \emptyset$ \\
Wähle $x_3$ \dots \\
Wähle $x_r \in X \setminus \bigcup\limits_{i=1}^{r-1}B_{\varepsilon}(x_i)$
\begin{align*}
	&\Rightarrow (X_r)_r \text{ Folge in }X 
\end{align*}	
	Für $n \neq m$ gilt: 
	$$d(x_n,x_m) \geq \varepsilon $$
	Dann ist $n>m$ \,,so ist nach Konstruktion  $x_n \notin B_{\varepsilon}(x_m)$ und 
	dann kann keine Teilfolge von $(x_r)_r$ eine Cauchyfolge sein.
 \[
 	\Rightarrow  (x_r)_r \text{ hat keine Konvergente Teilfolge} \]
	\[
		X \text{ nicht folgenkompakt}
	\]
\bewende


\newpage
\section{Differentialrechnung im $\mathbb{R}^n$} % (fold)
\label{sec:differentialrechnung_im_mathbb_r_n}

\subsection{Die Definition} % (fold)
\label{sub:die_definition}
\underline{Ziel}: $U \subseteq  \mathbb{R}^n$ offen, $f: U \to \mathbb{R}^m$ \\
Wann heißt $f$ differenzierbar? \\
Was ist die Ableitung von $f$? 
\subsubsection{Notation} % (fold)
\label{ssub:notation}
$U \subseteq \mathbb{R}^n, \, x \in \mathbb{R}^n$.
\[
	x+U = \set[x+y]{y \in U} \subseteq \mathbb{R}^n
\]
Analysis I: $n=m=1$, $U \subseteq \mathbb{R}$ offen, $f: U \to \mathbb{R} \,, x \in U$. $f$ diffbar $:\Leftrightarrow$ der Grenzwert
\[
	f'(x)= \lim_{h \to \infty} \frac{1}{h}(f(x+h)-f(x)) \in \mathbb{R}
\]
existiert. \\
\underline{Verallgemeinerung}: $\mathbb{R} \supseteq U \stackrel{f}{\to} \mathbb{R}^m$ ohne Problem. \\
n>1 problematisch, denn man kann nicht durch Elemente $h \in \mathbb{R}^n$ dividieren. \\
\underline{Umformulierung}:(aus Analysis I) \\
$U \subseteq \mathbb{R}$ offen, $x \in U$, $f: U \to \mathbb{R}$ 
\[
	A: -x+U \to \mathbb{R}  \qquad ( -x+U=\{h \in \mathbb{R}\,|\,h+x \in U \} )
\]
definiert durch:
\[
	A(h):= \begin{cases}
		\frac{1}{n}f(x+h)-f(x), &\text{ falls }h \neq 0\\
		f(x), &\text{ falls }h=0
		\end{cases}
\]
$\Rightarrow A \text{ stetig bei }0$
\[
	\forall\, h \in -x+U \qquad f(x+h)=f(x)+A(h)h
\]
Umgekehrt: Nimm an, es gibt Funktion $A:-x+U \to \mathbb{R}$ ,so dass 
\begin{enumerate}
	\item $A$ stetig bei $0$
	\item $f(x+h)=f(x)+A(h)h$
\end{enumerate}
$\Rightarrow f$ differenzierbar bei $x$, und $f'(x)=A(0)$
\minisec{Beweis:}
\[
	\frac{1}{h}(f(x+h)-f(x))=\frac{1}{h}(f(x)+A(h)h-f(x))= \frac{1}{h}A(h)h = A(h) \stackrel{h \to 0}{\longrightarrow} A(0)
\]
\subsubsection{Bemerkung} % (fold)
\label{ssub:bemerkung}
\begin{enumerate}
	\item für $n=m=1$ erhält man den bekannten Begriff aus Ana I 
	\item hier ist $\text{Mat}_{m,n}(\mathbb{R})$ isomorph zu $\mathbb{R}^{m \times n}$ mit irgendeiner Norm ausgestattet (welche, ist gleichgültig, weil alle Normen
	Äquivalent sind)
	\item $f= \begin{pmatrix}
		f_1 \\ \vdots \\ f_m
	\end{pmatrix}$ $f: U \to \mathbb{R}$ \\
	$f$ diffbar bei $x$ $\Leftrightarrow$ $f$ diffbar bei $x$ ($\forall\, i \in \{1,\dots,m\}$) \\
	und
	\[
		Df(x)= \begin{pmatrix}
			Df_1(x) \\
			\vdots \\
			Df_n(x)
		\end{pmatrix} \qquad Df(x) \in \Mat_{m,n}(\mathbb{R})
	\]
	\item Wir identifizieren Matrizen und lineare Abbildungen mittels Isomorphismen
	\begin{align*}
		\Mat_{m,n}(\mathbb{R}) &\rightarrow \Hom(\mathbb{R}^n;\mathbb{R}^m) \\
		 A &\mapsto (x \mapsto A \cdot x)
	\end{align*}
\end{enumerate}
% subsubsection bemerkung (end)
% subsubsection notation (end)
\subsubsection{Definition} % (fold)
\label{ssub:definition}
$U \subseteq \mathbb{R}^n$ offen, $f:U \to \mathbb{R}^m$ $x \in U$. f heißt {\bfseries differenzierbar} bei $x$, falls eine Funktion $A: -x+U \to \Mat_{m,n}(\mathbb{R})$ existiert, so dass gilt:
\begin{enumerate}
	\item $A$ stetig bei 0
	\item $\forall\, h \in -x+U:$
	\[
		\underset{\in \mathbb{R}^m}{\underbrace{f(x+h)}} = \underset{\in \mathbb{R}^m}{\underbrace{f(x)}} + \underset{\in \mathbb{R}^m}{\underbrace{\underset{
		\in \Mat_m,n(\mathbb{R})}{\underbrace{A(h)}}
		\underset{\in \mathbb{R}^n}{\underbrace{h}}}}
	\]
\end{enumerate}
Die Matrix $A(0) \in \Mat_{m,n}(\mathbb{R})$ heißt {\bfseries Ableitung} von $f$ in $x$ und wird mit
\[
	Df(x) \in \Mat_{m,n}(\mathbb{R})
\]
bezeichnet \\
\underline{Alternative Bezeichnung}:
'total differenzierbar' \\
\[
	Df(x)=\text{'Totales Diferential'}
\]
\subsubsection{Satz} % (fold)
\label{ssub:satz}
$U \subseteq \mathbb{R}^n$ offen, $x \in U$,$f: U \to \mathbb{R}^m$ sei differenzierbar. Sei $v \in \mathbb{R}^n$ \\
setze $g_v(t)=f(x+tv)$ mit 
\[
	g_v:\{t \in \mathbb{R} \,|\,x+tv \in U \} \to \mathbb{R}^m
\]
Dann gilt:
\[
	Df(x)v = \lim_{t \to 0} \frac{1}{t}(g_v(t)-g_v(0))
\]
\minisec{Beweis:}
Schreibe $f(x+h)=f(x)+A(h)h$. $A$ stetig bei $0$ und $A(0)=Df(x)$
\begin{align*}
	\Rightarrow \frac{1}{t}(g_v(t)-g_v(0))&= \frac{1}{t}(f(x+tv)-f(x))  \\
	&= \frac{1}{t}(f(x)+A(tv)tv-f(x)) \\
	&= \frac{1}{t}A(tv)tv \\
	&= \frac{1}{t}tA(tv)v = A(tv)v \longrightarrow A(0)v = Df(x)v \\
\end{align*}
\[
	t \mapsto A(tv)v \text{ ist stetig bei $0$}
\] \bewende
\minisec{Bemerkung:}
In Def II1.2 ist $A(h)$ {\bfseries nicht} eindeutig durch $f$ bestimmt. Aber $A(0)$ ist eindeutig bestimmt, nach II.1.3
% subsubsection satz (end)
\subsubsection{Definition} % (fold)
\label{ssub:definition}
$U,x,f,v$ wie in Satz II.1.3. Dann heißt $Df(x)v$ die {\bfseries Richtungsableitung} von $f$ in Richtung $v$.
\[
	D_vf(x)=Df(x)v
\] 
\minisec{Bemerkung:}
Ist $f$ diffbar bei $x$, so ist $v \mapsto g'_v(0)=D_vf(x)$ linear.
\minisec{Beispiel:}
\[
	f(x)=\skal{x}{x}=\norm{x}_2^2
\]
\[
	f(x+h)=\skal{x}{x}+\skal{h}{h}+2\skal{x}{h}=f(x)+\skal{h}{h}+2\skal{x}{h}
\]
$A(h) \in \text{Mat}_{n,m}(\mathbb{R})$
\[
	A(h)h= \skal{h}{h}+2 \skal{x}{h}   
\]
% subsubsection definition (end)
% subsubsection definition (end)
% subsection die_definition (end)
\subsubsection{Satz} % (fold)
\label{ssub:satz}\marginnote{die Numerierung ist hier nicht richtig}
$U \subseteq \mathbb{R}^n$ offen, $x \in U$, $f:U \to \mathbb{R}^m$. $C \in \Mat_{m,n}(\mathbb{R})$.\\ Es sind Äquivalent:
\begin{enumerate}
	\item $f$ ist differenzierbar bei $x$, und $Df(x)=C$
	\item Es gilt:
	\[
		\lim_{h \to 0} \frac{\norm{f(x+h)-f(x)-C \cdot h}}{\norm{h}} = 0
	\]
\end{enumerate}
2 besagt, dass $h \mapsto f(x)+C \cdot h$ eine gute Approximation an $h \mapsto f(x+h)$ ist) \\
Außerdem liefern Satz II.1.8 und II.1.5 ein 'Kochrezept' um die differenzierbarkeit von $f$ nachzuprüfen:
 \begin{enumerate}
 	\item Berechne $D_{v}f(x)$ für alle $v \in \mathbb{R}^n$
	\item Finde Matrix $C \in \Mat_{m,n}(\mathbb{R})$ mit ($v=D_vf(x),\forall\, v \in \mathbb{R}^n$)
	\item Prüfe \[
		\lim_{h \to 0} \frac{1}{\norm{h}}\norm{f(x+h)-f(x)-C \cdot h}=0
	\]
	\[
		\Rightarrow Df(x)=C
	\]
 \end{enumerate}
 
 \minisec{Beweis:} 
 \underline{$1 \Rightarrow 2$}: Sei $A: \,-x+U \to \Mat_{m,n}(\mathbb{R})$ mit
 \[
 	f(x+h)=f(x)+A(h)h \qquad (\forall\, h)
 \]
 $A$ stetig bei $0$ und $A(0)=C$
 \[
 	\Rightarrow \frac{\norm{f(x+h)-f(x)-C \cdot h}}{\norm{h}}= \frac{A(h)h-C \cdot h}{\norm{h}} \leq \underset{\substack{\to 0 \\ \text{für }h \to 0 \\ \text{($A$ stetig bei }0)}}{\underbrace{\norm{A(h)-C}}} \frac{\norm{h}}{\norm{h}}
 \]
% subsubsection satz (end)
\underline{$2 \Rightarrow 1$}: \\
\underline{Vorbemerkung}: $h \in \mathbb{R}^n$,$v \in \mathbb{R}^n$,$h \neq 0$. Dann ist durch 
\[
	y \mapsto \frac{\skal{h}{y}}{\norm{h}^2}v , \mathbb{R}^n \to \mathbb{R}^m 
\]
eine lineare Abbildung gegeben, welche mit $h \mapsto v$ \\
Setze $\varphi (h) = f(x+h)-f(x)-C \cdot h , \varphi: -x+U \to \mathbb{R}^m$ \\
Nach Voraussetzung 2) gilt:
\[
	\frac{\norm{\varphi(h)}}{\norm{h}} \stackrel{h \to 0}{\longrightarrow} 0 
\]
Definiere Funktion $A: -x+U \to \Mat_{m,n}(\mathbb{R})$, $A(0)=C \qquad h \neq 0$\\
Sei $A(h)$ die Matrix, welche die lineare Abbildung \[
	y \mapsto C \cdot y + \frac{\skal{h}{y}}{\norm{h}^2}\varphi(h)
\]
darstellt. wobei $\varphi(h)=A(h)y$
\begin{align*}
	\Rightarrow f(x)+A(h)h &= f(x)+C \cdot h + \frac{\skal{h}{h}}{\norm{h}^2}\varphi(h) \\
	&=f(x)+C \cdot h + \varphi(h) \\
	&=f(x+h)	\\
\end{align*}

Ferner ist $A$ stetig bei $0$ ,denn

\begin{align*}
	\norm{A(h)-A(0)} &= \sup\limits_{\norm{y}\leq 1} \norm{A(h)y- A(0)y} \\
			&= \sup\limits_{\norm{y}\leq 1} \norm{Cy + \frac{\skal{h}{y}}{\norm{h}^2} \varphi(h) - Cy } \\
			&= \sup\limits_{\norm{y}\leq 1} \frac{\abs{\skal{h}{y}}}{\norm{h}^2} \norm{\varphi(h)} \\
			&\leq \sup\limits_{\norm{y}\leq 1} \frac{\norm{h}\norm{y}}{\norm{h}^2} \norm{\varphi(h)} \\
			&=\frac{\norm{\varphi(h)}}{\norm{h}} \stackrel{h \to 0}{\longrightarrow} 0
\end{align*}
\[
	\Rightarrow f \text{ differenzierbar bei }x
\]
\[
	Df(x)=C
\]
\subsubsection{Definition} % (fold)
\label{ssub:definition}
$U \subseteq \mathbb{R}^n$ offen, $x \in U$, $f:  \to \mathbb{R}^m$. Die {\bfseries $i$-te partielle Ableitung von $f$} bei $x$ ist
\[
	D_if(x):= D_{e_i}f(x) \qquad (i=1,\dots,n)
\]
Man schreibt oft auch \[
	\diff{}{x_i}f(x):= D_if(x)
\]
Es gilt folgendes: \\
Wenn $f$ differenzierbar bei $x$ ist, so gilt: \\
Der $(i,j)$te Eintrag von $Df(x)$ ist \[
	(Df(x))_{i,j}=D_jf_i(x) \qquad i=1,\dots,m , j=1,\dots,n
\]
\[
	Df(x)= \begin{pmatrix}
		D_1f_1(x) & \dots & D_nf_1(x) \\
		\vdots & & \vdots \\
		D_1f_m(x) & \dots & D_nf_m(x)
	\end{pmatrix} \text{ (Jacobi-Matrix)}
\]
Wie berechnet man $D_jf_i(x)$?
\[
	D_jf_i(x):=  \diffd{}{t} \Big|_{t=0}f_i(x+te_j)= \diffd{}{t} \Big|_{t=0}f_i(x_1,\dots,x_{j-1},x_j+t,x_{j+1},\dots,x_n)
\]
\minisec{Beispiel:} $U= \set[(x_1,x_2) \in \mathbb{R}^2]{x_2 \not= 0}$ offen in $\mathbb{R}^2$ 
\[
	q(x_1,x_2)= \frac{x_1}{x_2}
\]
Berechnung der partiellen Ableitungen:
\[
	D_1q(x_1,x_2)= \diffd{}{t} \Big|_{t=0} q(x_1+t,x_2)= \diffd{}{t} \Big|_{t=0} \frac{x_1+t}{x_2} = \frac{1}{x_2} \left(= \diffd{}{x_2} \frac{x_1}{x_2} \right)
\]
\[
	D_2q(x_1,x_2)= \diffd{}{t}  \Big|_{t=0} q(x_1,x_2+t)= \diffd{}{t}  \Big|_{t=0} \frac{x_1}{x_2+t} = \frac{-x_1}{(x_2+t)^2}  \Big|_{t=0}^{} = - \frac{x_1}{x_2} 
	\left( = \diffd{}{x_2} \frac{x_1}{x_2}  \right)
\]
$\Rightarrow$ falls $q$ diffbar ist, so muss
\[
	Dq(x_1,x_2)= \left( \frac{1}{x_2}, -\frac{x_1}{x_2} \right)
\]
Ist $q$ diffbar? \\
$x=(x_1,x_2) \in U$ fest, $h=(h_1,h_2) \in -x+U$. Setze
\[
	A_x(h):= \left( \frac{1}{x_2+h_2} , \frac{-x_1}{x_2(x_2+h_2)}  \right)
\]
Dann gilt $q(x+h)=q(x)+A_x(h)h$, denn
\begin{align*}
	q(x)+A_x(h)h &= \frac{x_1}{x_2}+ \left( \frac{h_1}{x_2+h_2}-\frac{x_1h_1}{x_2(x_2+h_2)}  \right) \\
	&= \frac{x_1(x_2+h_2)+x_2h_1-x_1h_2}{x_2(x_2+h_2)} \\
	&= \frac{x_1x_2 + x_2h_1}{x_2(x_2+h_2)} \\
	&= \frac{x_1+h_1}{x_2+h_2} \\
	&= q(x+h)
\end{align*}
$h \mapsto A_x(h)$ stetig bei 0, denn

\[
	A_x(h)= \enbrace[\bigg]{ \underset{\to \frac{1}{x_2}}{\underbrace{\frac{1}{x_2+h_2}}} , \underset{\to \frac{-x_1}{x_2^2}}{\underbrace{-\frac{x_1}{x_2(x_2+h_2)}}}  }
\]
\minisec{Beispiel:}
$f: \mathbb{R}^n \to \mathbb{R}$ mit $f(x)=\skal{x}{x}$. \\
Wir haben gesehen, dass $f$ differenzierbar ist mit 
\[
	D_vf(x)= 2 \skal{x}{v} \qquad \forall\, x,v \in \mathbb{R}^n
\] 
\begin{align*}
	D_vf(x) &= \diff{}{t}  \Big|_{t=0}\skal{x+tv}{x+tv} \\
			&= \diff{}{t}  \Big|_{t=0} \left( \skal{x}{x}^2 + 2 \skal{x}{v} + t^2 \skal{v}{v} \right) \\
			&= 2 \skal{x}{v}
\end{align*}
Somit folgt, dass gilt: \[
	D_if(x)=2 \skal{x}{e_i} = 2 x_i
\]
\[
	Df(x)=(2x_1, \dots, 2x_n)
\]
Dieses hätte man auch sehen können, wenn man 
\[
	f(x)= \sum^{n}_{i=1}x_i^2
\]
betrachtet mit 
\[
	\diff{}{x_i}f(x)=2x_i
\]
% subsubsection definition (end)

\subsubsection{Satz} % (fold)
\label{ssub:satz}
$U \subseteq \mathbb{R}^n$ offen, $x \in U$ und $f:U \to \mathbb{R}^n$. Falls die partiellen Ableitungen $D_1f, \dots , D_nf$ auf ganz $U$ existieren, und stetig 
sind (also Funktionen $D_if: U \to \mathbb{R}^m$) dann ist $f$ differenzierbar in $x$

\minisec{Beweis:}
Weil $f$ genau dann differenzierbar ist, wenn alle $f_1 , \dots , f_n: U \to \mathbb{R} $ differenzierbar sind, genügt es den Fall $m=1$ zu betrachten.
\begin{align*}
	f(x+h) - f(x) &= f(x_1+h_1, \dots ,x_n+h_n) - f(x_1, \dots, x_n) \\
	&= \sum^{n}_{i=1} \left( f(x_1+h_1, \dots,x_i+h_i,x_{i+1} ,\dots,x_n) - f(x_1+h_1, \dots,x_{i-1}+h_{i-1},x_{i} ,\dots,x_n) \right)
\end{align*}
Nach dem Mittelwertsatz gibt es $\theta_1, \dots , \theta_n \in (0,1)$, so dass
\begin{gather*}
	f(x_1+h_1, \dots,x_{i-1}+h_{i-1},x_i+h_i,x_{i+1} ,\dots,x_n) - f(x_1+h_1, \dots,x_{i-1}+h_{i-1},x_i,x_{i+1} ,\dots,x_n) \\
	= D_if(x_1+h_1, \dots,x_{i-1}+h_{i-1},x_{i}+\theta_ih_i,x_{i+1} ,\dots,x_n)h_i 
\end{gather*}
% subsubsection satz (end)
\[
	\Rightarrow f(x+h)-f(x)= \sum^{n}_{i=1}D_if(x_1+h_1, \dots,x_{i-1}+h_{i-1},x_{i}+\theta_ih_i,x_{i+1} ,\dots,x_n)h_i 
\]
Setze $A(h) \in \Mat_{n,m}(\mathbb{R})$ mit
\[
	A(h):= D_if(x_1+h_1, \dots,x_{i-1}+h_{i-1},x_{i}+\theta_ih_i,x_{i+1} ,\dots,x_n)
\]
\[
	\Rightarrow f(x+h)=f(x)-A(h)h
\]
Weil $D_if$ stetig bei $x$ ist, gilt
\[
	A(h) \stackrel{h \to 0}{\longrightarrow} A(0)
\]
\bewende \\
Wir haben gesehen: 
\[
	f \text{ diffbar } \Rightarrow f \text{ partiell diffbar}
\]
\[
	f \text{ stetig partiell diffbar } \Rightarrow f \text{ diffbar}
\]
% section differentialrechnung_im_mathbb_r_n (end)
\subsubsection{Definition} % (fold)
\label{ssub:definition}
$f: U \to \mathbb{R}^m$ heißt {\bfseries stetig differenzierbar}, falls $f$ differenzierbar ist und falls $Df: U \to \text{Mat}_{n,m}(\mathbb{R})$ stetig ist.
% subsubsection definition (end)

\subsubsection{Kettenregel} % (fold)
\label{ssub:kettenregel}
$U \subseteq \mathbb{R}^n$ offen, $V \subseteq \mathbb{R}^m$ offen, $f:U \to V \mathbb{R}^m$, $g: V \to \mathbb{R}^p$, $x \in U$, $y=f(x) \in V$. Falls $f$ 
differenzierbar ist, dann ist $g \circ f: U \to \mathbb{R}^p$ differenzierbar, und es gilt:
\[
	D(g \circ f)(x) = Dg(f(x))Df(x)
\]

\minisec{Beweis:} Sei $A:\, -x+U \to \text{Mat}_{n,m}(\mathbb{R})$ und $f(x+h)=f(x)+A(h)\cdot h $ für alle $h$. $A$ stetig in $0$ \\
$B:\, -y+V \to \text{Mat}_{p,m}(\mathbb{R})$ und $g(y+k)=g(y)+B(k)\cdot k$ für alle $k \in -y+V$. $B$ stetig in $0$
\[
	\Rightarrow g \circ f(x+h)= g( \underset{=y}{\underbrace{f(x)}}+ \underset{=k}{\underbrace{A(h)h}}) = g(f(x))+B(A(h)h)A(h)h =: g(f(x))+C(h)h
\]
wobei 
\[
	C(h):= \underset{\text{Mat}_{p,n}(\mathbb{R})}{\underbrace{ \underset{\in \text{Mat}_{p,m}(\mathbb{R})}{\underbrace{B(A(h)h)}} \underset{ \, \, \in
	 \text{Mat}_{n,m}(\mathbb{R})}{\underbrace{A(h)}} }}
\]
\begin{align*}
	C(0)&=B(A(0)0)A(0)\\
	&= B(0)A(0) \\
	&= Dg(y)Df(x) \\
	&= Dg(f(x))Df(x)
\end{align*}
Außerden ist die Funktion 
\begin{align*}
	C: \,-x+U &\to \text{Mat}_{p,n}(\mathbb{R}) \\
	h &\mapsto C(h)
\end{align*}	
stetig. \bewende
% subsubsection kettenregel (end)

\minisec{Formulierung der Kettenregel mit partiellen Ableitungen:}
$f,g$ wie in II.1.14 
\[
	f= \begin{pmatrix}
		f_1 \\
		\vdots \\
		f_m
	\end{pmatrix} \qquad g= \begin{pmatrix}
		g_1 \\
		\vdots \\
		g_p
	\end{pmatrix}
\]
Dann gilt:
\[
	Dg(f(x))= \begin{pmatrix}
		D_1g_1(f(x)) & \dots & D_mg_1(f(x)) \\
		\vdots & & \vdots \\
		D_1g_p(f(x)) & \dots & D_mg_p(f(x))
	\end{pmatrix}
\]
\[
	Df(x)= \begin{pmatrix}
		D_1f_1(x) & \dots & D_nf_1(x) \\
		\vdots & & \vdots \\
		D_1f_m(x) & \dots & D_nf_m(x)
	\end{pmatrix}
\]
\[
	\Rightarrow D(g \circ f)(x)= 
	\begin{pmatrix}
			D_1g_1 & \dots & D_mg_1 \\
			\vdots & & \vdots \\
			D_1g_p & \dots & D_mg_p
		\end{pmatrix} \cdot \begin{pmatrix}
		D_1f_1(x) & \dots & D_nf_1(x) \\
		\vdots & & \vdots \\
		D_1f_m(x) & \dots & D_nf_m(x)
	\end{pmatrix}
\]
Der $(k,l)$-te Eintrag von $D(g \circ f)(x)$ ist 
\[
	\sum^{m}_{j=1}D_jg_k(f(x))D_lf_g(x)=D_lg_k(f(x)) = \sum^{m}_{j=1}\diff{g_k}{y_g}(f(x))\diff{f_g}{x_l}(x)
\]

\minisec{Beispiele für Kettenregel:}
$U \subseteq \mathbb{R}^n$ offen, $f: U \to \mathbb{R}$, $g: U \to \mathbb{R} \setminus \set{0}$ seien differenzierbar.
\begin{align*}
	\begin{pmatrix}
		f \\ g 
	\end{pmatrix} : U &\to  \mathbb{R} \times ( \mathbb{R} \setminus \set{0})  \\
	x &\mapsto \begin{pmatrix}
		f(x) \\ g(x)
	\end{pmatrix}
\end{align*}
ist differenzierbar, und
\[
	D(\begin{pmatrix}
		f \\ g
	\end{pmatrix}) = \begin{pmatrix}
		D_1f(x) & \dots & D_nf(x) \\
		D_1g(x) & \dots & D_ng(x)
	\end{pmatrix}
\]
Die Funktion 
\begin{align*}
	q: \mathbb{R} \times ( \mathbb{R} \setminus \set{0}) &\to  \mathbb{R} \\
	(y,z) &\mapsto \frac{y}{z} 
\end{align*}
ist differenzierbar mit
\[
	Dq(y,z)= \begin{pmatrix}
		\frac{1}{z} & -\frac{y}{z^2}
	\end{pmatrix}
\]
\begin{align*}
	\Rightarrow q \circ \begin{pmatrix}
		f \\ g
	\end{pmatrix} : U &\to \mathbb{R} \\
	x &\mapsto \frac{f(x)}{g(x)}
\end{align*} ist differenzierbar.
\begin{align*}
	D\left(q \circ \begin{pmatrix}
		f \\ g
	\end{pmatrix}\right)&= Dq\left(\begin{pmatrix}
		f(x) \\ g(x)
	\end{pmatrix}\right) D \left(\begin{pmatrix}
		f \\ g
	\end{pmatrix}\right)(x) \\ &= \begin{pmatrix}
		\frac{1}{g(x)} & -\frac{f(x)}{g(x)^2}
	\end{pmatrix} \cdot \begin{pmatrix}
		D_1f(x) & \dots & D_nf(x) \\
		D_1g(x) & \dots & D_ng(x)
	\end{pmatrix} \\
	&= \left( \frac{D_1f(x)}{g(x)} - \frac{f(x)D_1g(x)}{g(x)^2}, \dots, \frac{D_nf(x)}{g(x)} - \frac{f(x)D_nf(x)}{g(x)^2}  \right)\\
	&= \left( \frac{D_1f(x)g(x)-f(x)D_1g(x)}{g(x)^2}, \dots , \frac{D_nf(x)g(x)-f(x)D_ng(x)}{g(x)^2}  \right)
\end{align*}
Analog erhält man durch Betrachtung der Funktion $\mu : \mathbb{R}^2 \to  \mathbb{R}$ mit $\mu(y,z)=yz$, die {\bfseries Produktregel}

\minisec{Beispiel:}
$A \in \Mat_{m,n}(\mathbb{R})$, $F: \mathbb{R}^n \to \mathbb{R}^m$, $F(x)=A \cdot x$
\[
	\Rightarrow F \text{ ist differenzierbar und es gilt } DF(x)=A
\]
\minisec{Beweis:} 
\[
	F(x+h)=A(x+h)=Ax+Ah=F(x)+Ah
\] \bewende \\
$U \subseteq \mathbb{R}^p$ offen, $f: U \to \mathbb{R}^n$ differenzierbar. \\
$\stackrel{\text{Kettenregel}}{\Rightarrow}$ $F \circ f: U \to \mathbb{R}^m$ ist differenzierbar, und 
\[
	D(f \circ f)(x)= DF(f(x))Df(x) = A \cdot Df(x)
\]
wobei
\[
	F \circ f(x)=A \cdot f(x)
\]
Das heißt Differenzieren vertauscht mit linearen Abbildungen.

\minisec{Beispiel:}
$A= \begin{pmatrix}
	a_1 & a_2
\end{pmatrix} \in \Mat_{1,2}(\mathbb{R})$
\[
	f: U \to \mathbb{R}^2 \qquad f(x)= \begin{pmatrix}
		f_1(x) \\ f_2(x)
	\end{pmatrix}
\]
\[
	Af=a_1f_1(x)+a_2f_2(x)
\]
\[
	\Rightarrow D(Af)(x)= A Df(x) = (a_1D_1f_1(x) + a_2D_1f_2(x), \dots, a_1D_nf_1 + a_2D_nf_2(x)) = a_1Df_1 + a_2Df_2(x) 
\]
Dieses ist die Summenregel
\minisec{Bemerkung:}
$X,Y$ endliche-dimensionale $\mathbb{R}$-VR. Dann gibt es Normen auf $X$ und $Y$. \\
(z.B $(v_1, \dots , v_n)$ Basis von $X$ und $x = \sum^{n}_{i=1}x_iv_i \in X \,\,,x_i \in \mathbb{R}$ wobei $\norm{x}:= \sum^{n}_{i=1}\abs{x_i}$) \\
Je zwei Normen auf $X$ sind äquivalent (weil $\dim X < \infty$). \\
Also gibt es die Begriffe 'Konvergente Folge','Stetigkeit','offene Mengen', etc, unabhängig von der Wahl der Norm.

\subsubsection{Definition} % (fold)
\label{ssub:definition}
$X,Y$ endlich-dimensionale $\mathbb{R}$-VR, $U \subseteq X$ offen, $x \in U$,$f: U \to Y$ heißt {\bfseries differenzierbar} bei $x$, falls gilt: \\
es gibt eine Abbildung $A: \,-x+U \to \Hom(X,Y)$, so dass gilt:
\begin{enumerate}
	\item $f(x+h)=f(x)+\underset{\in \Hom(X,Y)}{\underbrace{A(h)}}(h) \in Y$ für alle $h \in -x+U$
	\item $A$ stetig bei $0$ (Diese Aussage ist wohldefiniert, weil $\Hom(X,Y)$ ein endlich-dimensionaler $\mathbb{R}$-VR ist)
	\[
		A(0)=Df(x)
	\]
\end{enumerate}
\subsection{Differentiation unter dem Integral} % (fold)
\label{sub:differentiation_unter_dem_integral}
Betrachte $f:[a,b] \times U \to \mathbb{R}$, $a,b \in \mathbb{R}$, $U \subseteq \mathbb{R}^n$ offen, $f$ stetig.
\begin{align*}
	F: U &\to \mathbb{R} \\
	F(x) &= \int\limits_{a}^{b}f(t,x) \,\mathrm{d}t
\end{align*}
\minisec{Frage:}
\begin{enumerate}
	\item ist $F$ stetig?
	\item falls $D_if$ existiert ($i=1,\dots,n$), existiert dann $D_iF$, und gilt
	\[
		D_iF(x) = \int\limits_{a}^{b}D_if(t,x) \,\mathrm{d}t ? 
	\]
\end{enumerate}
\subsubsection{Satz} % (fold)
\label{ssub:satz}
Seien $a,b \in \mathbb{R}$ und $U \subseteq \mathbb{R}^n$ offen. Sei $f:[a,b] \times U: \to \mathbb{R}$ stetig und setze $F(x):= \int\limits_{a}^{b}f(t,x) \,\mathrm{d}t$. Dann ist $F: U \to \mathbb{R}$ stetig.

\minisec{Beweis:} Sei $x \in U$ und sei $\varepsilon >0$. $U$ ist offen, es gibt also $r>0$ mit:
\[
	\bar B_r(x)= \set[y \in U]{\norm{x-y}\leq r} \subseteq U 
\]
$[a,b] \times \bar B_r(x)$ ist Kompakt (denn es ist abgeschlossen und beschränkte Teilmenge des $\mathbb{R}^{n+1}$)
\[
	f  \Big|_{[a,b] \times \bar B_r(x)}^{} : [a,b] \times \bar B_r(x) \to \mathbb{R}
\] ist gleichmäßig stetig!! \\
$\Rightarrow$ es gibt ein $\delta >0$, so dass gilt falls $y \in U$, $t \in [a,b]$ und $\norm{y-x}\leq \delta $ und $\norm{y-x}\leq r$
\[
	\Rightarrow \abs{f(t,y)-f(t,x)}<\varepsilon \qquad (\delta \leq r)
\] 
Falls $\norm{y-x}<\delta \leq r$, so gilt also:
\[
	\abs{F(x)-F(y)} = \abs{ \int\limits_{a}^{b}f(t,x)-f(t,y) \,\mathrm{d}t} \leq \int\limits_{a}^{b} \abs{f(t,x)-f(t,y)} \,\mathrm{d}t < \int\limits_{a}^{b} \varepsilon \,\mathrm{d}t = \varepsilon \abs{b-a}
\] \bewende
% subsubsection satz (end)
\subsubsection{Satz} % (fold)
\label{ssub:satz}
Seien $a,b \in \mathbb{R}$, $U \subseteq \mathbb{R}$ offen. Sei $f: [a,b]\times U \to \mathbb{R}$ stetig. Angenommen, es existert $D_2f(t,x)= \diff{}{x}f(t,x)$, und $D_2f:[a,b] \times U \to \mathbb{R}$ sei stetig. \\
Dann ist $F:U \to \mathbb{R}$ mit $F(x):= \int\limits_{a}^{b}f(t,x)\,\mathrm{d}t$ differenzierbar und es gilt:
\[
	F'(x)= \int\limits_{a}^{b}\diff{}{x}f(t,x) \,\mathrm{d}t
\]
\minisec{Beweis:} Seien $x,x+h \in U$
Dann gilt nach dem Hauptsatz der Differential- und Integralrechnung und weil $D_ef$ stetig:
\begin{align*}
	\frac{1}{h}(F(x+h)-F(x))- \int\limits_{a}^{b}D_2f(t,x) \,\mathrm{d}t &= \frac{1}{h}\int\limits_{a}^{b}( f(t,x+h)-f(t,x)-D_2f(t,x)h) \,\mathrm{d}t \\
	&\stackrel{*}{=} \frac{1}{h}\int\limits_{a}^{b}( \int\limits_{x}^{x+h}D_2f(t,x) \,\mathrm{d}y) - D_2f(t,y)h \,\mathrm{d}t \\
	&= \frac{1}{h}\int\limits_{a}^{b}(\int\limits_{x}^{x+h}(D_2f(t,x)-D_2f(t,y)) \,\mathrm{d}y) \,\mathrm{d}t \\
\end{align*} 
Nun sei $\varepsilon>0$. Wie im Beweis von II.2.1 gibt es ein $r>0$, $0 < \delta \leq r$, so dass gilt:
\[
	\abs{x-y}< \delta \,, t \in [a,b] \Rightarrow  \abs{D_2f(t,x)-D_2f(t,y)} < \varepsilon \qquad (D_2 \text{ stetig})
\]
Ist also $\abs{h}< \delta $,so folgt:
\begin{align*}
	\abs{\frac{1}{h}(F(x+h)-F(x))- \int\limits_{a}^{b}D_2f(t,x) \,\mathrm{d}t} &= \abs{\frac{1}{h}} \cdot
	\abs{\int\limits_{a}^{b}\left(\int\limits_{x}^{x+h}(D_2f(t,x)-D_2f(t,y)) \,\mathrm{d}y\right) \,\mathrm{d}t} \\
	& \leq \frac{1}{\abs{h}} \int\limits_{a}^{b}
	\bigg(\underset{<\abs{h}\varepsilon}{\underbrace{\int\limits_{x}^{x+h}\underset{<\varepsilon}{\underbrace{\abs{D_2f(t,x)-D_2f(t,y)}}}\,\mathrm{d}y}}\bigg)  \,\mathrm{d}t \\
	&< \frac{1}{\abs{h}}\abs{b-a}\cdot\abs{h} \\
	&= \varepsilon \abs{b-a} 
	\end{align*}
	\bewende

\subsubsection{Korollar} % (fold)
\label{ssub:korollar}
$U \subseteq \mathbb{R}^n$ offen. $f:[a,b] \times U \to \mathbb{R}$. $f$ sei stetig und $\diff{}{x_i}f(t,x), i=1,\dots,n$ existiere und sei stetig. Dann ist $F:U \to \mathbb{R}$ mit $F(x):=\int\limits_{a}^{b}f(t,x) \,\mathrm{d}t$ stetig differenzierbar, und es gilt:
\[
	\diff{}{x_i}F(x)= \int\limits_{a}^{b}\diff{}{x_i}f(t,x) \,\mathrm{d}t
\]
\minisec{Beweis:} Sei $x \in U$, und $I = \set[s \in \mathbb{R}]{x+se_i \in U} \subseteq \mathbb{R}^n$ offen \\
Sei $g:[a,b] \times I \to \mathbb{R}$ mit $f(t,s)=f(t,x+se_i)$. Dann ist $g$ stetig und $\diff{}{s}g$ existiert und ist stetig, weil $\diff{}{s}g(t,s)= \diff{}{x_i}f(t,x+se_i)$
\[
	\stackrel{II.2.3}{\Rightarrow} G(s)= \int\limits_{a}^{b}g(t,s) \,\mathrm{d}t \text{ ist differenzierbar, und} 
\]
\[
	G'(x)= \int\limits_{a}^{b}\diff{}{x_i}f(t,x+se_i) \,\mathrm{d}
\]
aber $G(s)= F(x+se_i)$, also $G'(0)= \diff{}{x_i}F(x)$. \\
$\Rightarrow$ alle partiellen Ableitungen von $F$ existieren, und es gilt:
\[
	\diff{}{x_i}F(t,x)= \int\limits_{a}^{b} \underset{\text{stetig!}}{\underbrace{\diff{}{x_i}f(t,x)}} \,\mathrm{d}t
\]
Damit folgt, dass $\diff{}{x_i}F(t,x)$ stetig ist. \\
$ \stackrel{II.1.12}{\Rightarrow}$ $F(t,x)$ ist differenzierbar, und $DF$ ist stetig.
\bewende
% subsubsection korollar (end)
% subsubsection satz (end)

\subsubsection{Satz von Schwarz} % (fold)
\label{ssub:satz_von_schwarz}
Sei $U \subseteq \mathbb{R}^n$ offen, $f: U \to \mathbb{R}$. $v,w \in \mathbb{R}^n$.\\
Annahmen: \begin{itemize}
	\item $f$ stetig
	\item $D_vf$,$D_wf$ existieren und seien stetig
	\item $D_vD_wf$ existiere und sei stetig
\end{itemize}
Dann gilt:
\[
	D_wD_vf \text{ existiert und es gilt } D_wD_vf= D_vD_wf
\]
\minisec{Beweis:} 
\underline{1.Spezialfall}: $n=2$, $U=B_r(0)$ mit $r>0$. $v=e_1$, $w=e_2$ \\
Nach dem Hauptsatz und weil $D_2f$ stetig ist gilt:
\[
	f(x,y)= f(x,0)+\int\limits_{0}^{y}D_2f(x,t) \,\mathrm{d}t
\]
\begin{align*}
	\Rightarrow D_1f(x,y)=D_1f(x,0)+D_1 \int\limits_{0}^{y}D_2f(x,t) \,\mathrm{d}t \\
	&=D_1f(x,0)+ \int\limits_{0}^{y}D_1D_2f(x,t) \,\mathrm{d}t
\end{align*}
Weil $D_1D_2$ stetig.
\[
	D_1f(x,y)=\underset{\substack{\text{hängt nicht} \\\text{ von $y$ ab} \\ \Rightarrow D_2D_1f(x,0)=0}}{\underbrace{D_1f(x,0)}}+ \underset{\substack{D_2 \left( \int\limits_{0}^{y}D_1D_2f(x,t) \,\mathrm{d}t \right) \\ \stackrel{HS}{=}D_1D_2f(x,y) }}{\underbrace{\int\limits_{0}^{y}D_1D_2f(x,t) \,\mathrm{d}t}}
\]
$\Rightarrow$ $D_2D_1f(x,y)$ existiert, und es gilt:
\[
	D_2D_1f(x,y)= D_1D_2f(x,y)
\]
\underline{2.Allgemeiner Fall}: Sei $x \in U$, $v,w \in \mathbb{R}^n$. \\
$0 \in \tilde U= \set[(s,t)\in \mathbb{R}^2]{x+st+tw \in U} \subseteq \mathbb{R}^2$ offen. \\
Sei $g:\tilde U \to \mathbb{R}$ mit $g(s,t)=f(x+sv+tw)$. Dann ist $g$ stetig und
\begin{align*}
	D_1g(s,t) &= D_vf(x+sv+tw) \\
	D_2g(s,t) &= D_wf(x+sv+tw) \\
	D_1D_2g(s,t) &= D_vD_wf(x+sv+tw)
\end{align*}
existieren und sind stetig. \\
Nach Schritt 1 existiert $D_2D_1g$ in $0$, und $$D_wD_vg(x)= D_2D_1g(0,0)=D_1D_2g(0,0)=D_vD_wf(x)$$
\bewende
% subsubsection satz_von_schwarz (end)
\subsubsection{Definition} % (fold)
\label{ssub:definition}
$U \subseteq \mathbb{R}^n$ offen, $f: U \to \mathbb{R}^m$. Dann heißt $f$ {\bfseries r-mal stetig differenzierbar} (oder $C^r$-Funktion), falls gilt: \\
 für alle $s=0,\dots,r$ und alle $v_1,\dots,v_s \in \mathbb{R}^m$ existiert die {\bfseries iterierte Richtungsableitung}
 \[
 	D_{v_1}D_{v_2} \cdot \dots \cdot D_{v_s}f = D_{v_1}(D_{v_2} \cdot \dots \cdot D_{v_s}) 
 \]
 auf ganz $U$, und ist stetig auf $U$.
 
 \minisec{Definition:} 
 $U \subseteq \mathbb{R}^n$ offen, $f: U \to \mathbb{R}$ eine $C^2-Funktion$. Die {\bfseries Hessematrix} von $f$ bei $x$ ist
 \[
 	Hf(x):= \begin{pmatrix}
 		D_1D_1f(x) & \dots & D_nD_1f(x) \\
		\vdots & & \vdots \\
		D_1D_nf(x) & \dots & D_nD_nf(x)
 	\end{pmatrix}
 \]
 $\Rightarrow$ $Hf(x)$ ist symmetrisch, denn
 \[
 	Hf(x)^T \stackrel{\text{Satz von Schwarz}}{=} Hf(x)
 \] 
 Ferner gilt: sind $v,w \in \mathbb{R}^n$, so ist
 \[
 	D_vD_wf(x)= \skal{Hf(x)v}{w}=\skal{v}{Hf(x)w}
 \]
 \minisec{Beweis:} \[
 	v= \sum^{n}_{i=1}V_ie_i \qquad w=\sum^{n}_{i=1}w_ie_i \qquad v_i,w_i \in \mathbb{R}
 \]
\begin{align*}
	\Rightarrow D_vD_wf(x)&= \sum^{n}_{i=1}v_iD_iD_wf(x) \\ &= \sum^{n}_{i=1}\sum^{n}_{j=1}v_iw_jD_iD_j(x) \\ &= \sum^{n}_{i=1}\sum^{n}_{j=1}v_iw_j(Hf(x))_{ji} \\
	&= \sum^{n}_{i=1}w_j(Hf(x)v)_j = \skal{Hf(x)v}{w}= \skal{v}{Hf(x)w}
\end{align*}
Die Identität
\[
	\skal{Hf(x)v}{w}=\skal{v}{Hf(x)w} 
\] folgt daraus, dass $Hf(x)$ symmetrisch ist, denn
\[
	\skal{Av}{w}=\skal{v}{A^Tw} \qquad A \in \Mat_{n,n}(\mathbb{R})
\]
\subsection{Taylorentwicklung} % (fold)
\label{sub:taylorentwicklung}
Der eindimensionale Fall \\
$(a,b) \subseteq \mathbb{R} \qquad f:(a,b) \to \mathbb{R}$
$f$ sei $m$-mal stetig differenzierbar, $p \in (a,b), h \in \mathbb{R}$ mit $p+h \in (a,b)$

\subsubsection{Lemma/Definition} % (fold)
\label{ssub:lemma_definition}
Sei $f:(a,b) \to \mathbb{R}$ $m$-mal stetig differenzierbar, $p \in (a,b)$. Dann gibt es genau ein Polynom $T_p^mf(h)$, so dass gilt
\begin{enumerate}[(i)]
	\item Der Graph von $T^m_pf$ ist $m$-mal stetig differenzierbar
	\item Es gilt
	\[
		\left( \diffd{}{h} \right)^k T_p^mf(0)= f^{(k)}(p) \qquad \forall\, k=0,\dots,m
	\]
\end{enumerate}
Dieses Polynom heißt $m$-tes {\bfseries Taylorpolynom} von $f$ bei $p$.
\minisec{Beweis:} Setze
\[
	T_p^mf(h)=\sum_{k=0}^{m}\frac{1}{k!}f^{(k)}(p)h^k
\]
Dann gilt $(i)$ und $(ii)$ $\Rightarrow$ Existenz
\underline{Eindeutigkeit}: Sei $q(h)$ ein Polynom mit $(i)$ und $(ii)$
\[
	q(h)= \sum_{k=0}^{m}c_kh^k \qquad c_k \in \mathbb{R}
\]
\[
	\begin{array}{ccc}
	q(0)=c_0 & \Rightarrow & c_0=f^{(0)}(p) \\
	q'(0)=c_1 & \Rightarrow & c_1=f^{(1)}(p) \\
	q''(0)=2c_2 & \Rightarrow & c_2=\frac{1}{2}f^{(2)}(p) \\
	\vdots & & \vdots \\
	q^{(l)}(0)=c_l & \Rightarrow & c_l=\frac{1}{l!}f^{(l)}(p) \\
	\end{array}
\] 
mit
\[
	q^(l)(h)= \sum_{k=l}^{m}c_kk(k-1) \cdot \dots \cdot (k-l)h^{k-l}
\]
\[
	\Rightarrow q(h)=\sum_{k=0}^{m}\frac{1}{k!}f^{(k)}(p)h^k
\] \bewende
\subsubsection{Taylorscher Satz} % (fold)
\label{taylor}
Sei $f:(a,b) \to \mathbb{R}$ $(m+1)$-mal stetig differenzierbar, $p,p+h \in (a,b)$. \\
Dann gilt
\[
	f(p+h)= \underset{= T^m_pf(h)}{\underbrace{\sum_{k=0}^{m}\frac{1}{k!}f^{(k)}(p)h^k}} 
	+ \underset{\text{Restglied}}{\underbrace{\frac{1}{m!}h^{m+1}\int_{0}^{1}(1-t)^mf^{(m+1)}(p+th) \,\mathrm{d}t}}
\]
\minisec{Beweis:} Induktion über $m$:
\begin{description}
	\item[$m=0$:] Hauptsatz
	\item[$m \mapsto m+1$:] siehe Übung A3B1. partielle Integration
	\[
	 (1-t)^m= - \frac{1}{m+1}\diffd{}{t}(1-t)^{m+1}
	\] 
\end{description}\bewende \\
Schreibe Restglied in einfacher Form
\subsubsection{Mittelwertsatz der Integralrechnung} % (fold)
\label{ssub:mittelwertsatz_der_integralrechnung}
Sei $[a,b] \subseteq \mathbb{R}$ und $f,g:[a,b] \to \mathbb{R}$ stetig. $g(x)\geq 0 \forall\, x \in [a,b]$. \\
Dann gilt
\[
	\exists\, \xi \in [a,b] \text{ mit } \int_{a}^{b}f(t)g(t) \,\mathrm{d}t = f(\xi)\int_{a}^{b}g(t) \,\mathrm{d}t
\]
\minisec{Beweis:}Aus dem Satz vom Minimum und Maximum folgt: $\exists\, x_{\min}.x_{\max} \in [a,b] \,\, \forall\, x \in [a,b]$ mit
\[
	f(x_{\min}):=c \leq f(x) \leq C =: f(x_{\max}) 
\]
\[
	cg(x) \leq f(x)g(x)\leq C g(x), \qquad g(x) \leq 0
\]
Dann folgt
\[
	c \int_{a}^{b}g(t) \,\mathrm{d}t \leq \int_{a}^{b}f(t)g(t) \,\mathrm{d}t \leq C \int_{a}^{b}g(t) \,\mathrm{d}t
\]
Wir definieren uns eine Hilfsfunktion 
\[
	F(x):=f(x)\int_{a}^{b}g(t) \,\mathrm{d}t
\] Dann ist $F$ stetig und es gilt
\[
	F(x_{\min})=c \int_{a}^{b}g(t) \,\mathrm{d}t \qquad, \qquad F(x_{\max})= C \int_{a}^{b}g(t) \,\mathrm{d}t 
\]
\[
	\Rightarrow F(x_{\min}) \leq \int_{a}^{b}f(t)g(t) \,\mathrm{d}t \leq F(x_{\max})
\]
Benutze nun den Zwischenwertsatz und erhalte
\[
	\exists\, \xi \in [x_{\min},x_{\max}] \text{ mit } F(\xi)= \int_{a}^{b}f(t)g(t) \,\mathrm{d}t
\]
\bewende

\subsubsection{Lagrange-Darstellung des Restgliedes} % (fold)
\label{ssub:lagrange_darstellung_des_restgliedes}
Seien $a,b,p,p+h \in \mathbb{R}$ und $f:[a,b]  \to \mathbb{R}$ wie in \ref{taylor}. 
Dann gilt
\[
	\exists\,\xi \in [0,1] \text{ mit } f(p+h)=T^m_pf(h)+ \frac{h^{(m+1)}}{(m+1)!}f^{(m+1)}(p+\xi h)
\]
\minisec{Beweis:} zu Zeigen ist:
\[
	\exists\,\xi \in [0,1] \text{ mit } \frac{1}{m!}h^{m+1} \int_{0}^{1}(1-t)^mf^{(m+1)(p+th)} \,\mathrm{d}t = \frac{1}{(m+1)!}f^{(m+1)}(p+th)
\]
Aus \ref{ssub:mittelwertsatz_der_integralrechnung} folgt
\begin{align*}
	\exists\,\xi \in [0,1] \text{ mit } \int_{0}^{1}(1-t)^m f^{(m+1)}(p+th) \,\mathrm{d}t 
	&= f^{(m+1)}(p+\xi h) \underset{\frac{1}{m+1}}{\underbrace{\int_{0}^{1}(1-t)^m \,\mathrm{d}t}} \\
	&= \frac{1}{m+1}f^{(m+1)}(p+\xi h)
\end{align*}
$\Rightarrow$ Vorfaktoren einsetzen
\bewende

Taylorpolynom
\[
	f(p+h)=\sum_{k=0}^{m}\frac{1}{k!}f^{(k)}(p)h^k + \frac{1}{(m+1)!}f^{(m+1)}(p+ \xi h)h^{(m+1)}
\]
\subsubsection{Bemerkung} % (fold)
\label{ssub:bemerkung}
Ist $f$ beliebig oft stetig differenzierbar, so heißt die Reihe
\[
	\sum^{\infty}_{k=0}\frac{1}{k!}f^{(k)}(p)h^k
\]
{\bfseries Taylorreihe} von $f$ bei p und ist eine Potenzreihe.
\begin{itemize}
	\item Die Taylorreihe von $f$ muss nicht konvergieren, und selbst wenn sie konvergiert braucht die Grenzfunktion nicht gleich $f$ zu sein.
	\minisec{Gegenbeispiel:}
	\[
		f:\mathbb{R}\to \mathbb{R}, \qquad f(x) = \begin{cases}
			e^{-\frac{1}{x^2}}, &\text{ falls }x \neq 0\\ 0, &\text{ falls }x=0 \end{cases}
	\] 
	$\Rightarrow $ $f$ beliebig oft stetig differenzierbar $f^{(k)}(0)=0 \forall\, k$
	$\Rightarrow $ Taylorreihe ist $0 \neq f$
	\item Ist $\sum_{k=0}^{\infty}a_kx^k$ eine Potenzreihe mit Konvergenzradius $R>0$, und $f(x)= \sum_{k=0}^{\infty}a_kx^k$. Dann gilt
	\[
		T_0^mf(h)=\sum_{k=0}^{m}a_kh^k \qquad \text{(Übungsaufgabe)}
	\]
\end{itemize}
Der mehrdimensionale Fall: \\
$ U \subseteq \mathbb{R}^n$ offen, $f:U \subseteq \mathbb{R}$ $(m+1)$-mal stetig differenzierbar, $x \in U$
% subsubsection bemerkung (end)
% subsubsection lagrange_darstellung_des_restgliedes (end)
\subsubsection{Definition} % (fold)
\label{ssub:definition}
$U \subseteq \mathbb{R}^n$, $x \in U$. Dann heißt $U$ {\bfseries sternförmig }bezüglich $x$, falls gilt:
\[
	y \in U,\,t \in [0,1] \Rightarrow (1-t)x+ty \in U
\]
Nun sei $p \in U \subseteq \mathbb{R}^n$ offen, $U$ steinförmig bzgl. $p$. Sei $f: U \to \mathbb{R}$ $(m+1)$-mal stetig differenzierbar. $v \in \mathbb{R}^n$ mit $p+v \in U$. Setze $g(t):=f(p+tv)$ mit
\[
	g: \set[t \in \mathbb{R}]{p+tv \in U} \to \mathbb{R}
\]
Dann gilt
\[
	g(1)= \sum_{k=0}^{m}\frac{1}{k!}g^{(k)}(0)+ \frac{1}{(m+1)!}g^{(m+1)}(\xi)
\]
\[
	\Rightarrow g^{(k)}(0)=\underset{k \text{-mal}}{\underbrace{D_v \dots D_vf(p)}}
\]
\minisec{Beweis:} 
mit Indutkion:
\begin{description}
	\item[$k=0$:] klar
	\item[$k-1 \mapsto k$] 
	\[
		g^{(k)}(t)=\diffd{}{t}g^{(k-1)}(t)= \diffd{}{t}(\underset{(k-1) \text{-mal}}{\underbrace{D_v \dots D_v}} f(x+tv))
		= \underset{k \text{-mal}}{\underbrace{D_v \dots D_v}} f(x+tv)
	\] 
\end{description}
\subsubsection{Taylorentwicklung in $\mathbb{R}^n$} % (fold)
\label{ssub:taylorentwicklung_in_rn}
Sei $U \subseteq \mathbb{R}^n$ sternförmig bzgl. $x \in U$, $ U \subseteq \mathbb{R}^n$ offen. $ f: U \to \mathbb{R} $ $(m+1)$-mal stetig differenzierbar und $v \in \mathbb{R}^n$,$x+v \in U$. \\
Dann gilt
\[
	\exists\,\xi \in [0,1] \text{ mit } f(x+v)
	=\sum_{k=0}^{m}\frac{1}{k!} \underset{k \text{-mal}}{\underbrace{D_v \dots D_vf(x)}}+ \frac{1}{(m+1)!}\underset{(m+1)
	 \text{-mal}}{\underbrace{D_v \dots D_v}}f(x+ \xi v)
\]
Interessanter Spezialfall $m=1$:
\begin{align*}
	f(x+v)&=f(x)+ D_vf(x)+ \frac{1}{2}D_vD_vf(x+\xi v) \\
	&= f(x) + Df(x) \cdot v + \frac{1}{2} \skal{Hf(x+ \xi v)v}{v}
\end{align*}
Dies ist die Taylorformel für $m=1$
% subsubsection taylorentwicklung_in_rn (end)
% subsubsection definition (end)
\subsection{Extremwertprobleme} % (fold)
\label{sub:extremwertprobleme}
$ U \subseteq \mathbb{R}^n$ offen , $ f: U \to \mathbb{R} $ differenzierbar
\subsubsection{Satz} % (fold)
\label{ssub:satz}
Falls $f$ bei $x \in U$ ein lokales Minimum oder Maximum hat, so gilt $Df(x)=0$. \\
Lokales Minimum bei $x$ bedeutet:
\[
	\exists\, \text{Umgebung }x \in V \subseteq U \text{, so dass} \forall\, y \in V: f(y) \geq f(x)
\]
(lokales Maximum analog mit $f(y)\leq f(x)$)
\minisec{Beweis:} 
Betrachte Funktion $g_v(t)=f(x+tv)$ ($x$ lokales Mininum/Maximum, $v \in \mathbb{R}^n$). \\
$g_v$ ist definiert auf einer offenen Umgebung $I \subseteq \mathbb{R}$ von $0$. $g_v$ ist differenzierbar (Kettenregel).
\\ $g_v$ hat lokales Minimum oder Maximum bei $0$
\[
	g_v(0)=0 \text{ (Analysis I)}
\]
aber $g_v(0)=D_vf(x)=Df(x)\cdot v$
\[
	\Rightarrow Df(x)=0
\] \bewende

% subsubsection satz (end)
\subsubsection{Definition} % (fold)
\label{ssub:definition}
$ f: U \to \mathbb{R} $  differenzierbar und $x \in U$. $x$ heißt {\bfseries kritischer Punkt} von $f$ genau dann, wenn \[
	Df(x)=0
\] 
Sei $f$ von nun an $2$-mal stetig differenzierbar. Sei $x \in U$ kritischer Punkt von $f$. \\
Wähle $r > 0$ mit 
\[
	\underset{\text{sternförmig}}{\underbrace{B_r(x)}}\subseteq U
\]
Dann gilt für alle $v \in U$ mit $\norm{v}<r$
\[
	f(x+v)=f(x)+\frac{1}{2}\skal{Hf(x+\xi v)v}{v} \text{ mit }\xi \in [0,1])
\]
\minisec{Bespiele für $\mathbb{R}^2 \to \mathbb{R}$:}
\begin{enumerate}
	\item $f_0(x,y)=x^2+y^2$
	\item $f_1(x,y)=x^2-y^2$
	\item $f_2(x,y)=-x^2-y^2$
\end{enumerate}
Alle drei haben den kritischen Punkt $(0,0)$ und es gilt
\[
	Hf_0= \begin{pmatrix}
		2 & 0 \\ 0 & 2 
	\end{pmatrix} ,\qquad Hf_1=\begin{pmatrix}
		2 & 0 \\ 0 & -2
	\end{pmatrix} ,\qquad Hf_2=\begin{pmatrix}
		-2 & 0 \\ 0 & -2
	\end{pmatrix}
\] $f_0$ hat lokales Maximum in $0$. \\
Graph von $f_0$ ist Parabolloid. Niveaulinien sind Kreise. $f_2$ analog und $f_2 = -f_0$. \\
$f_1$ hat kein lokales Minimum/Maximum
\minisec{Exkurs über symmetrische Matrizen:}
\subsubsection{Definition} % (fold)
\label{ssub:definition}
$A \in \Mat_{n,n}(\mathbb{R})$, $\lambda \in \mathbb{R}$ heißt {\bfseries Eigenwert} von $A$ genau dann, wenn
\[
	\exists\,x \in \mathbb{R}^n, x \neq 0 \text{ und } Ax=\lambda x
\]
Die {\bfseries Eigenvektoren} zum Eigenwert $\lambda$ sind die $x \in \mathbb{R}^n$ mit $Ax=\lambda x$
\[
	\Eig(A,\lambda)= \set[x \in \mathbb{R}^n]{Ax=\lambda x}
\]
Das {\bfseries Charakteristische Polynom} ist definiert durch
\[
	\chi_A(t)=\det(A-t \cdot I_n)
\]
und ist ein Polynom vom Grad $n$. \\
Die Eigenwerte von $A$ sind die Nullstellen des charakteristischen Polynoms. Das heißt es gibt höchstens $n$ verschiedene Eigenwerte.
% subsubsection definition (end)
% subsubsection definition (end)
\subsubsection{Spektralsatz für symmetrische Matrizen} % (fold)
\label{ssub:spektralsatz_fur_symmetrische_matrizen}
Sei $A \in \Mat_{n,n}(\mathbb{R})$ symmetrisch. Dann gibt es eine Orthogonalbasis von $\mathbb{R}^n$, welche nur aus Eigenvektoren von $A$  besteht. \\
Mit anderen Worten: \\
Es gibt $\lambda_1 \leq \dots \lambda_n \in \mathbb{R}$ und $v_1,\dots,v_n \in \mathbb{R}^n$ mit
\[
	\skal{v_i}{v_j}= \delta_{ij}, \qquad Av_i=\lambda_iv_i
\]
\minisec{Bemerkung:}
$B= (v_1, \dots, v_n)$ Basis
\[
	M_B^B(A)= \begin{pmatrix}
	\lambda & \dots & 0 \\ \vdots & & \vdots \\ 0 &\dots& \lambda_n
	\end{pmatrix}
\] Für den Beweis siehe Lineare Algebra
% subsubsection spektralsatz_fur_symmetrische_matrizen (end)
\subsubsection{Satz} % (fold)
\label{ssub:satz}
Sei $A \in S(n)$, $\lambda_1 \leq \dots \leq \lambda_n$ seien Eigenwerte von $A$. \\
Dann gilt
\begin{enumerate}
	\item $\forall\, x \in \mathbb{R}^n:\, \lambda_1 \norm{x}^2 \leq \skal{Ax}{x} \leq \lambda_n \norm{x}^2$
	\item $\set[\skal{Ax}{x}]{\norm{x}=1}=[\lambda_1,\lambda_n]$
\end{enumerate}
\minisec{Beweis:} 
\begin{enumerate}
	\item $(v_1,\dots,v_n)$ Orthogonalbasis aus Eigenvektoren mit $Av_i=\lambda_iv_i$. $x \in \mathbb{R}^n$ mit $x = \sum^{n}_{i=1}x_iv_i$. \\
	Dann gilt
	\begin{align*}
		\skal{Ax}{x}&= \sum^{n}_{i,j=1} \skal{Ax_iv_i}{x_jv_j} \\
		&= \sum^{n}_{i,j=1} x_ix_j \skal{Av_i}{v_j} \\
		&= \sum^{n}_{i,j=1} x_ix_j\lambda_i \skal{v_i}{v_j} \\
		&= \sum^{n}_{i,j=1} \lambda_ix_i^2 \leq \lambda_n \sum^{n}_{i=1}x_i^2 = \lambda_n \norm{x}^2 \\
		&\geq \lambda_1 \sum^{n}_{i=1}x_i^2 = \lambda_1 \norm{x}^2
	\end{align*}
	\item Betrachte
	\[
		f(t)=(A(\cos(t)v_n+\sin(t)v_1),(\cos(t)v_1+\sin(t)v_n))
	\] Dann ist $f$ stetig und $\norm{\cos(t)v_n+\sin(t)v_1}=1$
	\[
		f(0)=\skal{Av_1}{v_n}= \lambda_n \skal{v_n}{v_n}= \lambda_n
	\]
	\[
		f(\frac{\pi}{2})= \skal{Av_1}{v_1}=\lambda_1
	\]
	Mit Zwischenwertsatz gilt
	\[
		[\lambda_1,\lambda_n] \subseteq f\left(\left[0,\frac{\pi}{2}\right]\right)  
	\]
\end{enumerate}
\bewende
\subsubsection{Satz} % (fold)
\label{ssub:satz}
Seien $ U \subseteq \mathbb{R}^n$ offen , $ x \in U $, $ f: U \to \mathbb{R} $ zweimal stetig differenzierbar. Sei $x$ ein kritischer Punkt von $f$.
\begin{enumerate}
	\item Falls alle Eigenwerte von $Hf(x)$ positiv sind, so hat $f$ ein lokales Minimum bei $x$ (positiv definit)
	\item Falls $Hf(x)$ negativ definit, so hat $f$ ein lokales Maximum bei $x$
	\item Falls $f$ einen positiven Eigenwert hat, so hat $f$ bei $x$ kein lokales Minimum
	\item Falls $f$ einen negativen Eigenwert hat, so hat $f$ bei $x$ kein lokales Maximum
	\item In allen anderen Fällen kann durch Betrachtung von $Hf(x)$ keine Entscheidung getroffen werden.
\end{enumerate}
\minisec{Beweis:} 
$1) \Rightarrow 2$, $3 \Rightarrow 4$ gehe zu $-f$ über
\begin{description}
	\item[3)] Sei $\norm{v}=1$, $Hf(x)v=\lambda v, \, \lambda >0$, $t \in [-1,1]$. \\ 
		$\Rightarrow \exists\,\xi \in [0,1]$ mit  
	\begin{align*}
		f(x+tv)&=f(x)+ \frac{1}{2} \skal{Hf(x+t \xi)tv}{tv} \\
		&= f(x) + \underset{= \frac{t^2}{2}\lambda}{\underbrace{\frac{t^2}{2} \skal{Hf(x)v}{v}}} 
		+ \frac{t^2}{2} \underset{\text{klein, falls $\abs{t}$ klein}}{\underbrace{\skal{(Hf(x+\xi t)- t f(x))v}{v}}} \\
		&> f(x) \text{ für kleine $t$}	
	\end{align*}
	\item[1)] selbe Idee
\end{description}
% subsubsection satz (end)
% subsubsection satz (end)
% subsection extremwertprobleme (end)

\end{document}	